\documentclass[11pt]{article}
\usepackage[margin=1in]{geometry}
\usepackage{hyperref}
\usepackage{graphicx}
\graphicspath{{./}}
\usepackage{natbib}
\usepackage{doi}
\setcitestyle{numbers,square} %
\makeatletter
\renewcommand{\hyper@natlinkbreak}[2]{#1}
\makeatother
\usepackage{titlesec}

\usepackage{amsmath}
\usepackage{amsthm}
\usepackage{dsfont} %
\theoremstyle{plain}
\newtheorem{theorem}{Theorem}
\newtheorem{proposition}[theorem]{Proposition}
\newtheorem{lemma}[theorem]{Lemma}
\newtheorem{corollary}[theorem]{Corollary}
\newtheorem{conjecture}[theorem]{Conjecture}

\theoremstyle{definition}
\newtheorem{definition}{Definition}
\newtheorem{axiom}[definition]{Axiom}
\newtheorem{assumption}[definition]{Assumption}
\newtheorem{criterion}[definition]{Criterion}
\newtheorem{example}[definition]{Example}

\theoremstyle{remark}
\newtheorem{remark}{Remark}
\newtheorem{observation}[remark]{Observation}
\newtheorem{property}[remark]{Property}

\usepackage{etoolbox}
\BeforeBeginEnvironment{theorem}{\needspace{4\baselineskip}\addvspace{6pt}}
\BeforeBeginEnvironment{corollary}{\needspace{4\baselineskip}\addvspace{6pt}}
\BeforeBeginEnvironment{lemma}{\needspace{4\baselineskip}\addvspace{6pt}}
\BeforeBeginEnvironment{proposition}{\needspace{4\baselineskip}\addvspace{6pt}}
\BeforeBeginEnvironment{definition}{\needspace{4\baselineskip}\addvspace{6pt}}
\BeforeBeginEnvironment{axiom}{\needspace{4\baselineskip}\addvspace{6pt}}
\BeforeBeginEnvironment{assumption}{\needspace{4\baselineskip}\addvspace{6pt}}
\BeforeBeginEnvironment{remark}{\needspace{4\baselineskip}\addvspace{9pt}}

\usepackage{longtable}
\usepackage{booktabs}
\usepackage{array}
\usepackage{calc}
\usepackage{caption}
\newenvironment{tablenote}
  {\par\addvspace{2pt}\footnotesize\noindent}
  {\par\addvspace{4pt}}
\newcounter{none}

\usepackage{parskip}

\widowpenalty=10000
\clubpenalty=10000
\usepackage[bottom]{footmisc}
\predisplaypenalty=10000
\postdisplaypenalty=150
\raggedbottom
\setlength{\skip\footins}{1.5\baselineskip}
\usepackage{needspace}
\let\origsection\section
\renewcommand{\section}{\needspace{12\baselineskip}\origsection}
\let\origsubsection\subsection
\renewcommand{\subsection}{\needspace{7\baselineskip}\origsubsection}
\let\origsubsubsection\subsubsection
\renewcommand{\subsubsection}{\needspace{4\baselineskip}\origsubsubsection}
\let\origparagraphcmd\paragraph
\renewcommand{\paragraph}{\needspace{3\baselineskip}\origparagraphcmd}
\let\origsubparagraph\subparagraph
\renewcommand{\subparagraph}{\needspace{3\baselineskip}\origsubparagraph}

\interdisplaylinepenalty=10000

\setcounter{secnumdepth}{3}

\titleformat{\paragraph}
  {\normalfont\normalsize\bfseries}{\theparagraph}{1em}{}
\titlespacing*{\paragraph}
  {0pt}{3.25ex plus 1ex minus .2ex}{1.5ex plus .2ex}
\titleformat{\subparagraph}
  {\normalfont\normalsize\bfseries}{\thesubparagraph}{1em}{}
\titlespacing*{\subparagraph}
  {0pt}{3.25ex plus 1ex minus .2ex}{1.5ex plus .2ex}

\providecommand{\tightlist}{%
  \setlength{\itemsep}{0pt}\setlength{\parskip}{0pt}}

\newcommand{\pandocbounded}[1]{%
  \sbox0{#1}%
  \ifdim\wd0>\linewidth
    \resizebox{\linewidth}{!}{#1}%
  \else
    #1%
  \fi
}

\usepackage{amssymb}
\usepackage{newunicodechar}
\newunicodechar{✓}{$\checkmark$}
\newunicodechar{≠}{$\neq$}
\newunicodechar{≡}{$\equiv$}
\newunicodechar{≈}{$\approx$}
\newunicodechar{δ}{$\delta$}
\newunicodechar{−}{$-$}
\newunicodechar{⁰}{$^{0}$}
\begingroup\renewcommand\PackageWarning[2]{}%
\newunicodechar{¹}{$^{1}$}
\newunicodechar{²}{$^{2}$}
\newunicodechar{³}{$^{3}$}
\endgroup
\newunicodechar{⁴}{$^{4}$}
\newunicodechar{⁵}{$^{5}$}
\newunicodechar{⁶}{$^{6}$}
\newunicodechar{⁷}{$^{7}$}
\newunicodechar{⁸}{$^{8}$}
\newunicodechar{⁹}{$^{9}$}
\newunicodechar{⁻}{$^{-}$}
\let\mystRightarrow\rightarrow
\renewcommand{\rightarrow}{\ensuremath{\mystRightarrow}}

\emergencystretch=2em

\hypersetup{colorlinks=true, allcolors=blue}

\title{Thermodynamics of Cooperation: Cooperative Equilibrium}
\author{%
Keith Lostracco%
}
\date{}

\begin{document}
\maketitle

\begin{abstract}
In a repeated multi-agent game we derive the payoff matrix from
thermodynamic first principles rather than stipulating it, denominating
payoffs in energy. Resource energy, friction multipliers, exploitation
efficiency, and repair costs fix the entries; a friction state variable
encoding the cumulative memory of past violations enters the payoffs
endogenously. Mutual defection is then net-negative rather than merely
suboptimal whenever combined boundary damage exceeds a modest threshold
set by the repair multiplier, a regime we argue is physically typical.
Our central result concerns scale: non-zero background friction holds
the N-player public goods threshold strictly below 1 for every group
size, approaching 0.4 at the baseline parameters. Without friction the
threshold rises to a ceiling equal to the inverse of the cooperative
multiplier, recovering the canonical large-group free-rider collapse in
the marginal-surplus limit. In the two-player game the same mechanism
places the critical discount factor near 0.43, against the canonical
0.5, dropping to about 0.19 once network friction is included. Energy
denomination also selects among equilibria: universal cooperation is the
unique Nash equilibrium that is simultaneously Pareto-efficient and
welfare-maximizing, a selection criterion within the continuum of Folk
Theorem equilibria, which persists, so cooperation is singled out, not
made unique. Defection is evolutionarily unstable, self-extinguishing in
friction-coupled networks, and unprofitable even from a single
deviation, making punishment self-reinforcing without external
enforcement. The repeated-game machinery is classical; the contribution
is what physically calibrated payoffs imply for equilibrium selection
and scale.
\end{abstract}
\section{Introduction}\label{introduction}

The classical repeated Prisoner's Dilemma demonstrates that cooperation
can be sustained as a Nash equilibrium when agents are sufficiently
patient \citep{Fudenberg1986, Fudenberg1991}. The Folk Theorem
establishes this as a general principle: any feasible, individually
rational payoff vector is achievable in equilibrium for sufficiently
high discount factors. This generality, while powerful, renders the
theorem indeterminate: it sustains cooperation but also sustains
exploitation, with no internal mechanism to select among the continuum
of equilibria.

This paper constructs a repeated game in which the payoff matrix is not
stipulated but derived from physical parameters. \emph{Necessary
Constraints} \citep{TC-I}, \emph{Thermodynamic Friction} \citep{TC-IV},
and \emph{Strategic Entropy Injection} \citep{TC-II} establish the
physical infrastructure: boundary constraints as Lagrangian multipliers
on multi-agent optimization, thermodynamic friction as the energy cost
of constraint violation, and information-theoretic costs of strategic
signal corruption. We inherit those results and construct the two-player
and \(N\)-player games with payoffs denominated in energy units
(Joules). The specific results inherited from these papers are recalled,
without proof, in the Appendix.

The energy denomination imposes structural constraints on the admissible
payoff space that address the Folk Theorem's indeterminacy by supplying
a principled selection criterion among its equilibria. In particular:

\begin{enumerate}
\def\labelenumi{\arabic{enumi}.}
\tightlist
\item
  The payoff matrix is \emph{derived} rather than assumed, with the
  Prisoner's Dilemma ordering emerging as a theorem
  (Proposition~\ref{prop-physical-prisoners-dilemma}).
\item
  Mutual defection becomes net-negative (\(P < 0\)) rather than merely
  suboptimal once combined boundary damage exceeds \(1/\kappa\) (one
  half at \(\kappa = 2\)) --- the destructive regime --- because
  thermodynamic friction then makes conflict self-destructive; at milder
  damage it remains costly but net-positive. This two-player regime map
  is the \emph{mechanism} beneath the headline \(N\)-player result, not
  the load-bearing claim itself.
\item
  In the \(N\)-player public goods extension, cooperation is sustainable
  as a Nash equilibrium in arbitrarily large populations and uniquely
  maximizes system welfare there (\(SW^{CC} = N(\alpha-1)c\),
  Theorem~\ref{thm-cooperation-maximizes-welfare}) --- the paper's
  load-bearing conclusion, and one that is \emph{independent of the
  friction multiplier} \(\Phi\): the threshold \(\delta_N^*\) contains
  no \(\kappa\), \(\psi\), or \(\Phi\). Concretely, the frictionless
  cooperation threshold rises to a ceiling of \(1/\alpha\) as
  \(N \to \infty\), reproducing the canonical large-group collapse in
  the limit of marginally beneficial cooperation (\(\alpha \to 1\),
  ceiling \(\to 1\)). Non-zero background friction replaces that ceiling
  with \(c/(\alpha c + \phi_0 \bar{E})\), bounded strictly below \(1\)
  for every \(\alpha\) and every \(N\), approaching \(0.4\) at the
  baseline \(\alpha = 2\)
  (Corollary~\ref{cor-cooperation-strengthens-scale}).
\item
  The resulting two-player cooperation threshold
  \(\delta^* \approx 0.430\) is below the canonical \(0.5\), and drops
  further to \(\approx 0.185\) once network friction enters the
  punishment branch.
\item
  A friction state variable encoding the cumulative memory of past
  violations enters payoffs endogenously, making punishment
  self-reinforcing.
\end{enumerate}

The paper proceeds as follows: model and physical setup,
energy-denominated payoff model, two-player Prisoner's Dilemma ordering,
cooperation theorem in the repeated game, network friction effects,
\(N\)-player public goods extension, evolutionary stability, tit-for-tat
dynamics, welfare uniqueness, and discussion.

\subsection{Related Work}\label{related-work}

The game-theoretic results build on two established traditions:
classical and evolutionary game theory, and the theory of repeated
games. Our contribution is the energy-denominated payoff matrix derived
from the dissipative-system model of \citep{TC-I, TC-IV}, together with
the cooperation, stability, and welfare consequences that follow from
its structural properties.

\textbf{Classical and evolutionary game theory.} \citet{VonNeumann1944}
created the mathematics of strategic interaction. \citeauthor{Nash1950}
\citetext{\citeyear{Nash1950}; \citeyear{Nash1951}} proved that every
finite game has at least one equilibrium in mixed strategies.
\citet{MaynardSmith1973} fused game theory with Darwinian selection,
introducing the Evolutionarily Stable Strategy (ESS), a concept strictly
stronger than Nash equilibrium that identifies strategies resistant to
invasion by mutants. \citet{Hamilton1964a} explained altruism through
inclusive fitness (\(rB > C\)). \citet{Axelrod1984} demonstrated that
Tit-for-Tat wins iterated Prisoner's Dilemma tournaments, establishing
that cooperation can emerge among self-interested agents without central
authority. \citet{Taylor1978} and \citet{Weibull1995} developed
replicator dynamics connecting game-theoretic equilibria to evolutionary
population dynamics. This literature uses dimensionless utility
(``utils''), a quantity that does not correspond to any physical
observable and is not interpersonally comparable. Our framework replaces
utils with energy units, inheriting the physical constraints that energy
must satisfy (conservation, positivity of costs, irreversibility of
entropy increase).

\textbf{The Folk Theorem and equilibrium selection.}
\citet{Fudenberg1986} proved the definitive Folk Theorem: when agents
are sufficiently patient, any feasible and individually rational outcome
can be sustained as a subgame-perfect Nash equilibrium. This result
simultaneously supports cooperation and creates indeterminacy: the
theory predicts everything and therefore nothing.
\citet{OsborneRubinstein1994} provide a comprehensive treatment of
equilibrium concepts in repeated games. Our framework addresses the Folk
Theorem's indeterminacy with an equilibrium-selection criterion rather
than an elimination of equilibria: the continuum of equilibria persists,
but the derived Prisoner's Dilemma structure --- with mutual defection
net-negative in the destructive damage regime --- together with the
endogenous friction state variable and the \(N\)-player
welfare-maximization result single out universal cooperation as the
unique equilibrium that is also Pareto-optimal and welfare-maximizing.

\textbf{Statistical physics of cooperation.} A substantial physics
literature studies cooperation on lattices and networks
\citep[§5]{SzaboFath2007, PercEtAl2013, PercEtAl2017, SantosSantosPacheco2008, JusupEtAl2022}.
The physics enters through the dynamics, in Monte Carlo updating by a
Fermi function \citep{PercEtAl2013} whose noise parameter
\citet{SzaboFath2007} call a ``temperature'', and in results reported as
phase diagrams and universality classes
\citep{SzaboFath2007, PercEtAl2017}. In the public goods models central
to this literature the payoffs themselves remain stipulated abstract
utilities, with the cost of cooperation set to unity and the synergy
factor swept as a control parameter
\citep{PercEtAl2013, PercEtAl2017, SantosSantosPacheco2008}. The present
paper inverts that division of labor, leaving the evolutionary dynamics
conventional and fixing the payoff scale from physical parameters.
\citet{SzaboFath2007} state the disanalogy that motivates the inversion:
a physical theory typically has a single function, a Hamiltonian or
thermodynamic potential, whose extremum determines the state of the
whole system, whereas a game typically has as many functions to maximize
as it has interacting agents.

Two prior works connect payoffs themselves to thermodynamics.
\citet{AnttilaAnnila2011} identify the rate of entropy increase as a
universal payoff function; the payoff is a rate, non-negative almost
everywhere, and the signed per-transaction quantity, the free energy
consumed in a move, is a driving force rather than a payoff. They argue
qualitatively that a coalition is warranted only where each member gains
higher entropic status than by acting alone, and read the tragedy of the
commons as thermodynamically probable once exploitation depletes the
surroundings below the society's own energy content, redirecting flows
away from it. \citet{AdamiHintze2018} map the one-shot Prisoner's
Dilemma and the three-player public goods game onto Ising-type
Hamiltonians, setting iterated play aside; there the Hamiltonian is
constructed from the payoff matrix, and the temperature is the inverse
strength of selection.

\textbf{Energy-denominated and physically grounded payoffs.} The
ecological economics tradition \citep{Lotka1922, GeorgescuRoegen1971}
recognized that energy is the fundamental currency of biological and
economic activity, but did not build individual-level game theory on
that foundation. \citet{Hirshleifer1991} modeled conflict as
resource-consuming economic activity but used abstract utility. The
thermodynamic friction framework \citep{TC-IV} provides the missing
link: contest dissipation, boundary damage, and repair costs yield a
friction multiplier \(\Phi\) that denominates the Prisoner's Dilemma
payoff matrix entirely in energy units. The resulting payoff structure
(with mutual defection net-negative once \(\kappa\psi > 1\)) is derived
from physical parameters rather than stipulated, and its consequences
for equilibrium selection, evolutionary stability, and welfare
uniqueness are the central contributions of this work.

\section{Model and Physical Setup}\label{model-and-physical-setup}

\subsection{Energy-Denominated
Payoffs}\label{energy-denominated-payoffs}

We adopt the physical framework established in \emph{Necessary
Constraints} \citep{TC-I} and \emph{Thermodynamic Friction}
\citep{TC-IV}, summarized here to establish notation. Two agents contest
a resource of energy value \(E_R > 0\). Each selects from two
meta-strategies:

\begin{itemize}
\tightlist
\item
  \textbf{Cooperate (C):} Accept the constrained allocation (the
  variational equilibrium of \citep{TC-I}). Pay only the coordination
  cost \(c \ll E_R\).
\item
  \textbf{Defect (D):} Violate the constraint. Seize the resource by
  force or strategic entropy injection, triggering the conflict contest
  dynamics of \citep{TC-IV} and/or the communication corruption dynamics
  of \citep{TC-II}.
\end{itemize}

The damage fractions \(\psi^{\text{off}}, \psi^{\text{def}} \in [0,1]\)
(with combined boundary damage
\(\psi := \psi^{\text{off}} + \psi^{\text{def}}\)), repair multiplier
\(\kappa \geq 1\), and friction multiplier \(\Phi = 1 + \kappa\psi\)
retain their definitions from \citep{TC-IV}. The symbol \(\delta\)
denotes the discount factor exclusively; the boundary-damage fractions
carry the distinct glyph \(\psi\).

\subsection{The Friction State
Variable}\label{the-friction-state-variable}

Defection does not merely affect the immediate interaction; it injects
friction into the network, degrading all future cooperative
transactions. The friction dynamics, formalized in \citep{TC-IV},
involve a network friction coefficient \(\phi \geq 0\) that
parameterizes per-transaction overhead costs due to monitoring,
verification, and precautionary defense. Each boundary violation of
severity \(v\) injects \(\Delta\phi = \eta v\) into the friction
coefficient (where \(\eta > 0\) is the friction sensitivity), while in
the absence of new violations, friction decays at rate
\(\epsilon \in (0,1)\) per period.

The combined law of motion is:

\[\phi_{t+1} = (1 - \epsilon)\,\phi_t + \eta\,v_t \cdot \mathds{1}_{\{\text{violation at } t\}}\]

The cascading cost of a single violation amplifies at ratios of \(10^3\)
to \(10^4\) across the cooperative network \citep{TC-IV}: a violation of
severity \(v\) produces total cascading cost
\(C_{\text{cascade}} = \eta v M \bar{E} / \epsilon\) over the full
recovery period.

The friction state variable enters our analysis through the
network-adjusted payoffs
(§\ref{network-effects-the-friction-amplifier}). The key property is
persistence: violations today degrade payoffs for approximately
\(1/\epsilon\) future periods, coupling the repeated game to a dynamic
system absent from standard models.

\section{The Energy-Denominated Payoff
Model}\label{the-energy-denominated-payoff-model}

\subsection{Motivation: Why Energy Changes
Everything}\label{motivation-why-energy-changes-everything}

Standard game theory uses abstract utility values (often called
``points'' or ``utils''). The Prisoner's Dilemma, for example, assigns
payoffs like (3, 3) for mutual cooperation and (1, 1) for mutual
defection, without specifying what these numbers represent or where they
come from.

Our framework replaces abstract utils with \textbf{thermodynamic energy}
(Joules, Calories, or normalized energy units). This is not merely a
relabeling: it imports the physical constraints that energy must
satisfy:

\begin{enumerate}
\def\labelenumi{\arabic{enumi}.}
\tightlist
\item
  \textbf{Conservation:} Energy is neither created nor destroyed.
  Conflict does not generate new energy; it dissipates existing energy
  as waste heat \citep{TC-IV}.
\item
  \textbf{Positivity of costs:} All physical actions require energy
  expenditure (\(W = \int F \cdot dx \geq 0\)). Conflict investments,
  boundary damage, and repair are non-negative real costs, physically
  determined rather than stipulated.
\item
  \textbf{Diminishing returns:} The concavity of utility in resources
  (the regularity assumption Assumption 1 \citep{TC-I}, part (b)) means
  that marginal gains decrease while marginal losses increase, a
  thermodynamic asymmetry that breaks the symmetry of abstract payoffs.
\item
  \textbf{Network coupling:} An agent's payoff depends not only on the
  current interaction but on the network friction state \(\phi\), which
  persists across interactions \citep{TC-IV}.
\end{enumerate}

These physical constraints \textbf{structurally determine} the payoff
matrices. The cooperative equilibrium is not an assumption fed into the
model; it is a \textbf{consequence} of the physics.

\subsection{Strategy Space}\label{strategy-space}

Each agent selects from two meta-strategies in any pairwise interaction
over a contested resource:

\begin{definition}[Cooperative and Defection Strategies]
\label{def-cooperative-defection-strategies}

Each agent selects one of two meta-strategies per interaction:

\begin{itemize}
\item \textbf{Cooperate (C):} Respect the boundary constraint (Definition 3 \citep{TC-I}). Accept the constrained allocation $\mathbf{x}_i^*$ determined by the variational equilibrium (Theorem 7 \citep{TC-I}). Incur only the trade/coordination cost $c > 0$ (the energy cost of negotiation, contract enforcement, or social coordination).
\item \textbf{Defect (D):} Violate the boundary constraint. Attempt to seize the resource by force or deception. This triggers the conflict contest (Definition 3 \citep{TC-IV}) and/or deception channel (Definition 3 \citep{TC-II}), incurring the associated costs.
\end{itemize}

\end{definition}

\textbf{Interpretation:} ``Cooperate'' means operating within the
constraint system. ``Defect'' means violating constraints. The
distinction is energetic: whether to expend resources on coordination or
on conflict.

\subsection{The Energy Variables}\label{the-energy-variables}

\begin{longtable}[]{@{}
  >{\raggedright\arraybackslash}p{(\linewidth - 4\tabcolsep) * \real{0.1111}}
  >{\raggedright\arraybackslash}p{(\linewidth - 4\tabcolsep) * \real{0.5556}}
  >{\raggedright\arraybackslash}p{(\linewidth - 4\tabcolsep) * \real{0.3333}}@{}}
\caption{Energy variables for a pairwise interaction over resource
\(E_R > 0\).}\tabularnewline
\toprule\noalign{}
\begin{minipage}[b]{\linewidth}\raggedright
Symbol
\end{minipage} & \begin{minipage}[b]{\linewidth}\raggedright
Meaning
\end{minipage} & \begin{minipage}[b]{\linewidth}\raggedright
Source
\end{minipage} \\
\midrule\noalign{}
\endfirsthead
\toprule\noalign{}
\begin{minipage}[b]{\linewidth}\raggedright
Symbol
\end{minipage} & \begin{minipage}[b]{\linewidth}\raggedright
Meaning
\end{minipage} & \begin{minipage}[b]{\linewidth}\raggedright
Source
\end{minipage} \\
\midrule\noalign{}
\endhead
\bottomrule\noalign{}
\endlastfoot
\(E_R\) & Energy value of the contested resource & Physical
environment \\
\(c\) & Coordination cost (trade negotiation, contract) & Cooperative
overhead; \(c \ll E_R\) \\
\(e^*\) & Nash equilibrium contest investment per agent & Proposition 1
\citep{TC-IV}: \(e^* = E_R/4\) (symmetric) \\
\(\Phi\) & Total friction multiplier & \citep{TC-IV}:
\(\Phi = 1 + \kappa\psi\),
\(\psi = \psi^{\text{off}} + \psi^{\text{def}}\) \\
\(\phi_0\) & Current network friction coefficient & Definition 8
\citep{TC-IV} \\
\(\Delta\phi\) & Friction injection from defection & Definition 9
\citep{TC-IV}: \(\Delta\phi = \eta \cdot v\) \\
\(\bar{E}\) & Average cooperative transaction value & Network
parameter \\
\(M\) & Number of transactions per period & Network parameter \\
\end{longtable}

\subsection{Admissibility: What Energy Denomination
Forbids}\label{admissibility-what-energy-denomination-forbids}

\begin{remark}[Admissibility of Energy-Denominated Payoff Matrices]
\label{rmk-admissibility}

The energy denomination restricts the class of admissible payoff matrices to those satisfying three constraints:

(a) \textbf{Resource-bounded gains.} All payoffs are bounded above by the resource endowment: $T, R, P, S \leq E_R$. No player can extract more energy than the resource contains. Losses are not symmetrically bounded: conflict can destroy more energy than the resource provides \citep{TC-IV}, e.g., $P < -E_R$ once $\Phi > 6$, though in every regime studied in this paper ($\Phi \leq 3.4$) all payoffs also satisfy $|T|, |R|, |P|, |S| \leq E_R$.

(b) \textbf{Strictly positive conflict costs.} Offensive and defensive actions require $e^* > 0$ energy investment. Costless aggression is physically inadmissible.

(c) \textbf{Persistent state coupling.} Payoffs in period $t$ depend on the network friction state $\phi_t$, which is a function of all prior interactions. Payoffs are not independently and identically distributed across periods.

Any payoff matrix violating (a)–(c) is physically inadmissible under energy denomination. Constraints (a)–(b) restrict the payoff space; constraint (c) restricts the game form.

A further property follows from (a)–(b) combined with the friction multiplier (Theorem 6 \citep{TC-IV}), conditional on the magnitude of the damage fractions:

\textbf{Conditional consequence: Net-negative mutual defection.} The mutual defection payoff satisfies $P = \frac{E_R}{2}(1 - \Phi/2) < 0$ exactly when $\Phi > 2$, which is equivalent to $\kappa(\psi^{\text{off}} + \psi^{\text{def}}) > 1$, i.e., to the combined boundary damage $\psi = \psi^{\text{off}} + \psi^{\text{def}}$ exceeding $1/\kappa$ (one half at $\kappa = 2$). The Second Law supplies the ingredients of this condition but not its magnitude: boundary damage exists on both sides ($\psi > 0$, since costless conflict is excluded by (b)), and repair is thermodynamically irreversible since repair requires thermodynamic work exceeding the damage energy ($\kappa > 1$, by the Second Law), but positivity and irreversibility alone do not force $\kappa\psi > 1$. Whether $\Phi > 2$ holds is therefore a regime condition on damage magnitudes rather than a derived necessity: it identifies the \textbf{destructive regime} ($\psi > 1/\kappa$), in which mutual conflict is self-destructive, and separates it from the mild regime ($\psi < 1/\kappa$) in which mutual defection is costly but net-positive. The two-agent crossover $\kappa\psi > 1$ is the same threshold established in \citep{TC-IV}; for three or more contestants it falls to $\Phi > N/(N-1)$, equivalently $\kappa\psi > 1/(N-1)$ — a threshold that shrinks as contestants are added and vanishes only in the large-$N$ limit, where any positive damage suffices. See §\ref{construction} for the full derivation.

In particular, the standard Prisoner's Dilemma with $P > 0$ is the mild-friction regime of this model, arising when $\Phi < 2$, i.e., $\kappa\psi < 1$. At $\kappa = 2$ this requires combined damage $\psi < 1/2$, which includes the baseline calibration ($\psi = 0.4$, giving $P = +5$) alongside lower-stakes disputes. Once combined damage exceeds $1/\kappa$, $\Phi > 2$ and $P < 0$.

\textbf{What is excluded:} Payoff specifications with unbounded gains, costless aggression, or memoryless (i.i.d.) interaction histories are inadmissible. Positive mutual-defection payoffs ($P > 0$) are not excluded by fiat: they arise in the mild regime $\Phi < 2$ (including the baseline) and become impossible once combined damage exceeds $1/\kappa$ (conditional consequence above). These exclusions are non-trivial: they eliminate entire families of payoff matrices that are routinely studied in abstract game theory but cannot arise when payoffs are denominated in conserved physical quantities subject to thermodynamic friction.
\end{remark}

\section{The Two-Player Payoff
Matrix}\label{the-two-player-payoff-matrix}

\subsection{Construction}\label{construction}

Consider agents \(A\) and \(B\) interacting over resource \(E_R\). Using
the energy variables defined above, with the equal-split convention
under cooperation:

\textbf{Case (C, C): Mutual Cooperation:}

Both agents respect the constraint. The resource is split according to
the variational equilibrium (equal split for symmetric agents). Each
pays the coordination cost \(c\).

\[\pi_A^{CC} = \pi_B^{CC} = \frac{E_R}{2} - c\]

\textbf{Case (D, C): Unilateral Defection (A defects, B cooperates):}

Agent \(A\) attacks the cooperating agent \(B\). Since \(B\) is not
investing in contest (expecting cooperation), \(A\) seizes the resource
but still incurs: - The metabolic cost of the aggressive action:
\(e_A^{\text{attack}}\) (energy of mounting the violation) - Repair of
offensive self-damage:
\(\kappa \cdot \psi^{\text{off}} \cdot e_A^{\text{attack}}\) (the
destroyed embodied energy \(\psi^{\text{off}} e_A^{\text{attack}}\) is
re-supplied by this repair term, not charged separately)

Agent \(B\), caught unprepared, suffers: - No resource share: \(B\)
receives \(0\) from the resource (forgoing the \(E_R/2\) it would have
obtained under cooperation) - Repair of defensive boundary damage:
\(\kappa \cdot \psi^{\text{def}} \cdot e_A^{\text{attack}}\) (again the
damage \(\psi^{\text{def}} e_A^{\text{attack}}\) is re-supplied by
repair, not counted twice)

However, the defector against a cooperator faces much lower resistance.
We model the attack cost against an undefended boundary as
\(e_A^{\text{attack}} = \theta \cdot E_R\) where \(\theta \in (0, 1)\)
is the \textbf{exploitation efficiency}, the fraction of the resource's
value the defector must spend to seize it from an unprepared cooperator.
Typically \(\theta \ll 1/2\) (attacking an undefended target is cheap).

\[\pi_A^{DC} = E_R - \theta E_R \left[1 + \kappa\psi^{\text{off}}\right] = E_R\left[1 - \theta(1 + \kappa\psi^{\text{off}})\right]\]

\[\pi_B^{DC} = -\theta E_R \kappa\psi^{\text{def}}\]

The victim \(B\) receives nothing from the resource and additionally
suffers boundary damage and repair costs from the attack.

\textbf{Case (C, D): Symmetric mirror of (D, C):}

\[\pi_A^{CD} = -\theta E_R \kappa\psi^{\text{def}}\]

\[\pi_B^{CD} = E_R\left[1 - \theta(1 + \kappa\psi^{\text{off}})\right]\]

\textbf{Case (D, D): Mutual Defection:}

Both agents fight. This is the Tullock contest from the thermodynamic
friction derivation (Proposition 1 \citep{TC-IV}). Each invests
\(e^* = E_R/4\), wins with probability 1/2, and pays the repair cost of
the boundary damage inflicted through both offense and defense (the
destroyed embodied energy is re-supplied by repair, so it is not charged
a second time as a standalone loss).

Each agent's expected payoff is the winning prize minus all costs:

\[\pi_A^{DD} = \pi_B^{DD} = \frac{E_R}{2} - e^* - \kappa(\psi^{\text{off}} + \psi^{\text{def}}) e^*\]

\[= \frac{E_R}{2} - \left[1 + \kappa(\psi^{\text{off}} + \psi^{\text{def}})\right] e^* = \frac{E_R}{2} - \Phi \cdot \frac{E_R}{4}\]

\[= \frac{E_R}{2}\left(1 - \frac{\Phi}{2}\right)\]

\subsection{The Payoff Matrix
(Single-Round)}\label{the-payoff-matrix-single-round}

Assembling into standard normal form, define the shorthand payoffs:

\[T = \pi^{DC}_{\text{defector}} = E_R\left[1 - \theta(1 + \kappa\psi^{\text{off}})\right] \quad \text{(Temptation)}\]

\[R = \pi^{CC} = \frac{E_R}{2} - c \quad \text{(Reward for cooperation)}\]

\[P = \pi^{DD} = \frac{E_R}{2}\left(1 - \frac{\Phi}{2}\right) \quad \text{(Punishment for mutual defection)}\]

\[S = \pi^{CD}_{\text{sucker}} = -\theta E_R \kappa\psi^{\text{def}} \quad \text{(Sucker's payoff)}\]

The normal-form game with row player \(A\) and column player \(B\) is
then:

\[\begin{array}{c|c|c}
 & B: C & B: D \\ \hline
A: C & (R,\; R) & (S,\; T) \\ \hline
A: D & (T,\; S) & (P,\; P) \\
\end{array}\]

\subsection{Parameterization: The Physical Prisoner's
Dilemma}\label{parameterization-the-physical-prisoners-dilemma}

The classic Prisoner's Dilemma requires \(T > R > P > S\). Let us verify
each inequality under physical constraints:

\textbf{Claim: \(T > R\) (temptation exceeds reward) when exploitation
is cheap.}

\[T > R \iff E_R[1 - \theta(1 + \kappa\psi^{\text{off}})] > \frac{E_R}{2} - c\]

\[\iff \frac{E_R}{2} + c > \theta E_R(1 + \kappa\psi^{\text{off}})\]

\[\iff \theta < \frac{1/2 + c/E_R}{1 + \kappa\psi^{\text{off}}}\]

For \(c \ll E_R\) and \(\psi^{\text{off}} = 0.15\), \(\kappa = 2\):
threshold \(\approx 0.5/1.3 \approx 0.385\). So \(T > R\) when the
exploitation cost \(\theta < 0.385\), i.e., when stealing from a
cooperator is less than 38.5\% of the resource value. This is the
standard condition: defection is \emph{tempting} when exploiting a
cooperator is cheap.

\textbf{Claim: \(R > P\) (cooperation beats mutual defection).}

\[R > P \iff \frac{E_R}{2} - c > \frac{E_R}{2}\left(1 - \frac{\Phi}{2}\right) = \frac{E_R}{2} - \frac{\Phi E_R}{4}\]

\[\iff \frac{\Phi E_R}{4} > c\]

Since \(\Phi \geq 1\) and \(c \ll E_R\), this holds for any non-trivial
friction multiplier. Even with zero boundary damage (\(\Phi = 1\)),
mutual cooperation beats mutual defection as long as the coordination
cost is less than \(E_R/4\), i.e., less than the per-agent contest
dissipation.

\textbf{Claim: \(P > S\) (mutual defection beats being the sucker) when
friction is not too extreme.}

\[P > S \iff \frac{E_R}{2}\left(1 - \frac{\Phi}{2}\right) > -\theta E_R\kappa\psi^{\text{def}}\]

\[\iff \frac{1}{2} - \frac{\Phi}{4} + \theta\kappa\psi^{\text{def}} > 0 \iff \Phi < 2 + 4\theta\kappa\psi^{\text{def}} \equiv \Phi_{\text{PD}}\]

For the baseline parameters (\(\theta = 0.15\), \(\kappa = 2\),
\(\psi^{\text{def}} = 0.25\)):
\(\Phi_{\text{PD}} = 2 + 4(0.15)(2)(0.25) = 2.3\). Since the baseline
\(\Phi = 1.8 < 2.3\), the PD ordering \(P > S\) holds. For
\(\Phi > \Phi_{\text{PD}}\), the game transitions to a \textbf{Chicken}
(Hawk-Dove) structure with \(T > R > S > P\): mutual defection becomes
\emph{worse} than being the sucker, because the contest dissipation at
high friction overwhelms the sucker's one-sided damage. In the Chicken
regime the grim-trigger threshold changes form: the deviator's best
response to perpetual punishment is then \(C\) rather than \(D\), so the
threshold is \((T-R)/(T-S)\) rather than \((T-R)/(T-P)\) (see
Theorem~\ref{thm-cooperation-theorem} and the sensitivity-table note
below). Numerically the transition remains \emph{favorable} for
cooperation, because the sucker payoff \(S\) worsens with friction and
enlarges \(T - S\), and even the one-shot game admits cooperative
mixed-strategy Nash equilibria under the Chicken structure.

\textbf{Claim: \(S < 0\) (the sucker always loses).}

\[S = -\theta E_R \kappa\psi^{\text{def}} < 0\]

since all parameters are positive. The sucker not only fails to obtain
the resource but suffers boundary damage from the defector's attack.

\medskip

\begin{proposition}[Physical Prisoner's Dilemma]
\label{prop-physical-prisoners-dilemma}

Under the energy-based payoff model with physical constraints ($E_R > 0$, $c \ll E_R$, $\theta < (1/2 + c/E_R)/(1 + \kappa\psi^{\text{off}})$, $1 < \Phi < \Phi_{\text{PD}} = 2 + 4\theta\kappa\psi^{\text{def}}$), the single-round game has the structure of a Prisoner's Dilemma: $T > R > P > S$.

In the single-round game, defection is the dominant strategy for each agent. The unique Nash Equilibrium is $(D, D)$ with payoffs $(P, P)$, even though $(C, C)$ with payoffs $(R, R)$ is Pareto-superior.
\end{proposition}

\begin{proof}
The inequalities $T > R > P > S$ are verified above under the stated parameter conditions. In a simultaneous single-round game:

\begin{itemize}
\item If $B$ plays $C$: $A$ prefers $D$ ($T > R$).
\item If $B$ plays $D$: $A$ prefers $D$ ($P > S$).
\end{itemize}

So $D$ strictly dominates $C$ for each player. The unique NE is $(D, D)$. But $R > P$ means both agents would be strictly better off at $(C, C)$ than at $(D, D)$, establishing the Pareto inefficiency of the one-shot equilibrium.
\end{proof}

\textbf{Physical interpretation:} Each agent's individually rational
choice (defect regardless of the opponent's action) leads to a
collectively irrational outcome. The inefficiency is stronger than in
the standard abstract PD: because \(P\) becomes net-negative once
\(\kappa\psi > 1\) (the destructive regime), mutual defection can yield
not merely a smaller gain but an absolute energy loss. Both agents would
prefer the cooperative outcome \((R, R)\), yet neither can unilaterally
switch to \(C\) without exposure to the sucker's payoff \(S < P\) (and
\(S < 0\) in every regime). The one-shot game therefore traps both
agents in the mutual-defection equilibrium --- net-negative in the
destructive regime --- that the repeated game
(§\ref{the-repeated-game-cooperation-as-equilibrium}) resolves.

\subsection{Numerical Illustration: The Baseline
Game}\label{numerical-illustration-the-baseline-game}

\textbf{Parameters:} \(E_R = 100\), \(c = 2\), \(\theta = 0.15\),
\(\psi^{\text{off}} = 0.15\), \(\psi^{\text{def}} = 0.25\),
\(\kappa = 2\).

Computed payoffs:

\[T = 100[1 - 0.15(1 + 2 \times 0.15)] = 100[1 - 0.15 \times 1.3] = 100 \times 0.805 = 80.5\]

\[R = 50 - 2 = 48\]

\[\Phi = 1 + 2(0.15 + 0.25) = 1.8\]

\[P = 50(1 - 1.8/2) = 50(1 - 0.9) = 50 \times 0.1 = 5\]

\[S = -0.15 \times 100 \times 2 \times 0.25 = -7.5\]

\begin{longtable}[]{@{}lll@{}}
\caption{Baseline single-round payoff matrix (energy
units).}\tabularnewline
\toprule\noalign{}
& \(B: C\) & \(B: D\) \\
\midrule\noalign{}
\endfirsthead
\toprule\noalign{}
& \(B: C\) & \(B: D\) \\
\midrule\noalign{}
\endhead
\bottomrule\noalign{}
\endlastfoot
\(A: C\) & \textbf{(48, 48)} & (−7.5, 80.5) \\
\(A: D\) & (80.5, −7.5) & (5, 5) \\
\end{longtable}

\textbf{Observations:}

\begin{enumerate}
\def\labelenumi{\arabic{enumi}.}
\tightlist
\item
  \textbf{Mutual defection is costly but net-positive at the mild
  baseline} (\(P = +5\)): the contest dissipation consumes most of the
  resource's value, leaving each agent only a small positive residual.
  Net-negativity (\(P < 0\)) requires the destructive regime
  \(\kappa\psi > 1\) (\(\psi > 1/2\) at \(\kappa = 2\)); see the worked
  example below and Theorem 6 \citep{TC-IV}.
\item
  \textbf{Mutual cooperation yields a large surplus} (\(R = 48\) each,
  system total = 96).
\item
  \textbf{Temptation is high} (\(T = 80.5\)), but the victim suffers
  (\(S = -7.5\)).
\item
  \textbf{The cooperation premium} is \(R - P = 48 - 5 = 43\) per agent,
  the energy saved by cooperative behavior compared to the conflict
  equilibrium.
\end{enumerate}

\subsection{Numerical Illustration: The Destructive
Regime}\label{numerical-illustration-the-destructive-regime}

The baseline above sits in the mild regime (\(\kappa\psi = 0.8 < 1\)),
where mutual defection is costly but net-positive. To exhibit a
genuinely net-negative conflict equilibrium we raise the damage
fractions into the destructive regime.

\textbf{Parameters:} \(E_R = 100\), \(c = 2\), \(\theta = 0.15\),
\(\psi^{\text{off}} = 0.25\), \(\psi^{\text{def}} = 0.35\) (so
\(\psi = 0.6 > 1/\kappa = 0.5\)), \(\kappa = 2\).

\[\Phi = 1 + 2(0.25 + 0.35) = 2.2\]

\[T = 100[1 - 0.15(1 + 2 \times 0.25)] = 100[1 - 0.15 \times 1.5] = 77.5\]

\[R = 50 - 2 = 48, \qquad P = 50(1 - 2.2/2) = -5, \qquad S = -0.15 \times 100 \times 2 \times 0.35 = -10.5\]

\begin{longtable}[]{@{}lll@{}}
\caption{Destructive-regime single-round payoff matrix (energy
units).}\tabularnewline
\toprule\noalign{}
& \(B: C\) & \(B: D\) \\
\midrule\noalign{}
\endfirsthead
\toprule\noalign{}
& \(B: C\) & \(B: D\) \\
\midrule\noalign{}
\endhead
\bottomrule\noalign{}
\endlastfoot
\(A: C\) & \textbf{(48, 48)} & (−10.5, 77.5) \\
\(A: D\) & (77.5, −10.5) & (−5, −5) \\
\end{longtable}

Here mutual defection is net-negative (\(P = -5\)): contest dissipation
plus repair now exceed the resource value, consistent with Theorem 6
\citep{TC-IV} (\(\Phi = 2.2 > 2\)). The Prisoner's Dilemma ordering is
preserved
(\(\Phi = 2.2 < \Phi_{\text{PD}} = 2 + 4(0.15)(2)(0.35) = 2.42\)), and
the cooperation threshold falls to
\(\delta^* = (77.5 - 48)/(77.5 - (-5)) = 29.5/82.5 = 0.358\), below the
baseline \(0.430\): harsher conflict makes cooperation easier to
sustain. This destructive regime, not the mild baseline, is where the
``mutual defection is net-negative'' reading of the model applies.

\section{The Repeated Game: Cooperation as
Equilibrium}\label{the-repeated-game-cooperation-as-equilibrium}

\subsection{Why Repetition Matters}\label{why-repetition-matters}

The single-round Prisoner's Dilemma has defection as the unique NE. But
entities in the real world (biological organisms, economic agents,
nations) interact \textbf{repeatedly} over indefinite timeframes. The
theory of repeated games \citep{Fudenberg1986} shows that sustained
cooperation can emerge as an equilibrium in the repeated version of the
game.

In the standard (abstract-utility) PD, this result is qualitative:
cooperation \emph{can} be sustained for sufficiently patient agents, but
the critical patience threshold \(\delta^*\) depends on payoff values
that are free parameters. In our physical PD, the threshold is
\emph{structurally constrained}. Because mutual defection is heavily
discounted by contest dissipation --- and net-negative once
\(\kappa\psi > 1\) (the destructive regime) --- the punishment for
defection is severe relative to the cooperative reward, and the
resulting \(\delta^*\) is low. The repeated-game analysis that follows
is standard in its mechanism (trigger strategies, folk theorem) but
novel in its conclusion: the physics, once the damage magnitudes are
calibrated, sets the threshold rather than the modeler.

The key parameter is the \textbf{discount factor} \(\delta \in (0, 1)\),
representing how much agents value future payoffs relative to present
ones. In our thermodynamic framework, \(\delta\) has a natural physical
interpretation:

\begin{definition}[Discount Factor]
\label{def-discount-factor}

The discount factor $\delta$ for an agent with per-period survival probability $\sigma \in (0, 1]$ is:

$$\delta = \sigma \cdot e^{-r}$$

where $r \geq 0$ is the time-preference rate. For entities optimizing over long horizons (high $\sigma$, low $r$), $\delta \to 1$. For entities with short horizons (low $\sigma$ or high $r$), $\delta \to 0$.
\end{definition}

\textbf{Physical interpretation:} \(\delta\) close to 0 means the agent
has a short planning horizon and discounts future payoffs heavily.
\(\delta\) close to 1 means the agent expects many future interactions
and weights them nearly equally with the present; this is formalized in
§\ref{extension-to-n-player-games} below.

\subsection{Trigger Strategies: Grim and
Tit-for-Tat}\label{trigger-strategies-grim-and-tit-for-tat}

The simplest mechanism to sustain cooperation in repeated games is a
\textbf{trigger strategy}, a rule that starts cooperating and retaliates
against defection.

\begin{definition}[Grim Trigger]
\label{def-grim-trigger}

Agent $i$ plays $C$ in every period as long as no agent has ever played $D$. If any defection occurs, agent $i$ plays $D$ forever.
\end{definition}

\begin{definition}[Tit-for-Tat (TFT)]
\label{def-tit-for-tat}

Agent $i$ plays $C$ in period 1. In period $t > 1$, agent $i$ plays whatever the opponent played in period $t - 1$.
\end{definition}

Both strategies embody a simple enforcement mechanism: cooperation is
maintained as long as agents observe cooperation; defection triggers
retaliation. The threat of retaliation, sustained loss of cooperative
payoffs, creates a cost to defection that extends beyond the single
interaction.

\subsection{The Cooperation Condition (Grim
Trigger)}\label{the-cooperation-condition-grim-trigger}

Under the grim trigger strategy, the following holds:

\begin{theorem}[Cooperation as Nash Equilibrium: The Cooperation Theorem]
\label{thm-cooperation-theorem}

In the infinitely repeated energy-based game with discount factor $\delta$ and the Prisoner's Dilemma ordering ($P > S$), the cooperative outcome $(C, C)$ in every period is sustainable as a Nash Equilibrium (under grim trigger) if and only if:

$$\delta \geq \frac{T - R}{T - P} = \delta^*$$

where $T$, $R$, $P$ are the temptation, reward, and punishment payoffs. In the Chicken regime ($S > P$, arising when $\Phi > \Phi_{\text{PD}}$), the deviator's best response to perpetual punishment is $C$ rather than $D$, and the corresponding threshold is $\delta^* = (T - R)/(T - S)$. In both regimes the critical discount factor $\delta^*$ has the physical interpretation: cooperation is sustainable when agents sufficiently value future interactions.
\end{theorem}

\begin{proof}
Consider agent $A$'s incentive to deviate from the cooperative profile $(C, C, C, \ldots)$.

\textbf{Payoff from continued cooperation:}

$$V_A^{\text{coop}} = R + \delta R + \delta^2 R + \cdots = \frac{R}{1 - \delta}$$

\textbf{Payoff from one-shot defection (then punished forever):}

The deviating agent gets $T$ in the current period, after which the opponent plays $D$ forever. In the PD regime ($P > S$) the deviator's best per-period reply to perpetual $D$ is $D$, so the punishment stream pays $P$:

$$V_A^{\text{deviate}} = T + \delta P + \delta^2 P + \cdots = T + \frac{\delta P}{1 - \delta}$$

Cooperation is incentive-compatible when $V_A^{\text{coop}} \geq V_A^{\text{deviate}}$:

$$\frac{R}{1 - \delta} \geq T + \frac{\delta P}{1 - \delta}$$

$$R \geq (1 - \delta) T + \delta P$$

$$R - P \geq (1 - \delta)(T - P)$$

$$\frac{R - P}{T - P} \geq 1 - \delta$$

$$\delta \geq 1 - \frac{R - P}{T - P} = \frac{T - R}{T - P}$$

Since both agents face the same calculation (the game is symmetric), $(C, C)$ sustained by grim trigger is a NE whenever $\delta \geq \delta^*$. In the Chicken regime ($S > P$) the deviator's best per-period reply to perpetual $D$ is instead $C$, earning $S > P$; replacing $P$ with $S$ throughout the same calculation yields the threshold $(T - R)/(T - S)$.
\end{proof}

\subsection{\texorpdfstring{Computing \(\delta^*\) Under Physical
Constraints}{Computing \textbackslash delta\^{}* Under Physical Constraints}}\label{computing-delta-star}

Using the energy-based payoffs from
§\ref{parameterization-the-physical-prisoners-dilemma}:

\[\delta^* = \frac{T - R}{T - P}\]

\textbf{Baseline numerical example}
(§\ref{numerical-illustration-the-baseline-game} parameters):

\[\delta^* = \frac{80.5 - 48}{80.5 - 5} = \frac{32.5}{75.5} = 0.4305\]

\begin{corollary}[Low Cooperation Threshold]
\label{cor-low-cooperation-threshold}

Under the physical payoff model, the critical discount factor satisfies $\delta^* < 0.5$ if and only if $T + P < 2R$. At the baseline parameters this exact condition holds directly ($T + P = 85.5 < 96 = 2R$), giving $\delta^* = 0.430 < 0.5$: agents need only value the future at more than 43%

(i) $\Phi > 2$ (net-negative mutual defection, $P < 0$), and

(ii) $c < \frac{\theta E_R}{2}\bigl(1 + \kappa\psi^{\text{off}}\bigr)$ (the coordination cost is less than half the exploitation overhead, equivalently $R > T/2$)

hold; conditions (i)–(ii) are sufficient but not necessary, and at the mild baseline (i) fails while $\delta^* < 0.5$ still holds via the exact condition.
\end{corollary}

\begin{proof}
For the exact condition: since $T - P > 0$,

$$\delta^* = \frac{T - R}{T - P} < \frac{1}{2} \iff 2(T - R) < T - P \iff T + P < 2R$$

For sufficiency of (i)–(ii): when $P < 0$ (which occurs exactly when $\Phi > 2$):

$$T - P = T + |P| > T$$

$$\delta^* = \frac{T - R}{T - P} < \frac{T - R}{T} = 1 - \frac{R}{T}$$

Condition (ii) is $R > T/2$, i.e., $1 - R/T < 0.5$: substituting the payoff expressions from the physical parameterization (§\ref{parameterization-the-physical-prisoners-dilemma}), $R = E_R/2 - c$ and $T = E_R[1 - \theta(1 + \kappa\psi^{\text{off}})]$, the condition $R > T/2$ reduces to exactly (ii). Together, (i) and (ii) give $T + P < T < 2R$.

Condition (ii) is a genuine magnitude restriction, not a consequence of the PD ordering: $c \ll E_R$ does not imply it when exploitation is very cheap (small $\theta$). For example, $(E_R, c, \theta, \kappa, \psi^{\text{off}}, \psi^{\text{def}}) = (100, 2, 0.006, 2, 0.02, 0.49)$ satisfies the full PD ordering with $\Phi = 2.02 > 2$ and $P = -0.5 < 0$, yet the right side of (ii) is $0.312 < 2 = c$ and $\delta^* = 51.376/99.876 = 0.514 > 0.5$. The claim $\delta^* < 0.5$ therefore requires (ii) and does not follow from $\Phi > 2$ alone.

For the baseline values ($c = 2$, $\theta = 0.15$, $E_R = 100$, $\kappa = 2$, $\psi^{\text{off}} = 0.15$, $\psi^{\text{def}} = 0.25$), mutual defection is net-positive ($P = 5 > 0$, since $\Phi = 1.8 < 2$), so sufficient condition (i) does not apply; the exact condition certifies the threshold directly, $T + P = 80.5 + 5 = 85.5 < 96 = 2R$, giving $\delta^* = 0.430 < 0.5$. In the destructive regime ($\psi > 1/\kappa$, e.g. $\psi^{\text{off}} = 0.25$, $\psi^{\text{def}} = 0.35$) condition (i) holds and the sufficient route also applies: the right side of (ii) is $0.15 \times 100 \times 1.5 / 2 = 11.25 \gg 2 = c$, so $\delta^* < 0.5$ there as well.
\end{proof}

\textbf{Physical interpretation:} Because mutual defection under
physical energy payoffs is heavily eroded by contest dissipation --- and
net-negative once \(\kappa\psi > 1\) --- the punishment for defection is
severe relative to the cooperative reward: agents earn far less, and in
the destructive regime actively \emph{lose} energy. Provided
exploitation is not so cheap that the temptation dwarfs the reward (the
exact condition \(T + P < 2R\), sufficient condition (ii)), this makes
cooperation easy to sustain: even agents with modest foresight
(\(\delta > 0.43\) at baseline) find it in their interest to cooperate.
Compare this with the abstract PD (e.g., \(T=5, R=3, P=1, S=0\)) where
\(\delta^* = 0.5\). At realistic magnitudes, physical constraints make
cooperation \textbf{more} accessible, not less.

\subsection{\texorpdfstring{Sensitivity Analysis: \(\delta^*\) Across
Parameter
Regimes}{Sensitivity Analysis: \textbackslash delta\^{}* Across Parameter Regimes}}\label{sensitivity-analysis-delta-star}

\begin{longtable}[]{@{}
  >{\raggedright\arraybackslash}p{(\linewidth - 16\tabcolsep) * \real{0.2609}}
  >{\raggedright\arraybackslash}p{(\linewidth - 16\tabcolsep) * \real{0.0435}}
  >{\raggedright\arraybackslash}p{(\linewidth - 16\tabcolsep) * \real{0.0870}}
  >{\raggedright\arraybackslash}p{(\linewidth - 16\tabcolsep) * \real{0.0870}}
  >{\raggedright\arraybackslash}p{(\linewidth - 16\tabcolsep) * \real{0.0870}}
  >{\raggedright\arraybackslash}p{(\linewidth - 16\tabcolsep) * \real{0.0870}}
  >{\raggedright\arraybackslash}p{(\linewidth - 16\tabcolsep) * \real{0.0870}}
  >{\raggedright\arraybackslash}p{(\linewidth - 16\tabcolsep) * \real{0.1304}}
  >{\raggedright\arraybackslash}p{(\linewidth - 16\tabcolsep) * \real{0.1304}}@{}}
\caption{Cooperation threshold \(\delta^*\) across friction and
exploitation-cost regimes.}\tabularnewline
\toprule\noalign{}
\begin{minipage}[b]{\linewidth}\raggedright
Regime
\end{minipage} & \begin{minipage}[b]{\linewidth}\raggedright
\(\kappa\)
\end{minipage} & \begin{minipage}[b]{\linewidth}\raggedright
\(\theta\)
\end{minipage} & \begin{minipage}[b]{\linewidth}\raggedright
\(\Phi\)
\end{minipage} & \begin{minipage}[b]{\linewidth}\raggedright
\(T\)
\end{minipage} & \begin{minipage}[b]{\linewidth}\raggedright
\(R\)
\end{minipage} & \begin{minipage}[b]{\linewidth}\raggedright
\(P\)
\end{minipage} & \begin{minipage}[b]{\linewidth}\raggedright
\(S\)
\end{minipage} & \begin{minipage}[b]{\linewidth}\raggedright
\(\delta^*\)
\end{minipage} \\
\midrule\noalign{}
\endfirsthead
\toprule\noalign{}
\begin{minipage}[b]{\linewidth}\raggedright
Regime
\end{minipage} & \begin{minipage}[b]{\linewidth}\raggedright
\(\kappa\)
\end{minipage} & \begin{minipage}[b]{\linewidth}\raggedright
\(\theta\)
\end{minipage} & \begin{minipage}[b]{\linewidth}\raggedright
\(\Phi\)
\end{minipage} & \begin{minipage}[b]{\linewidth}\raggedright
\(T\)
\end{minipage} & \begin{minipage}[b]{\linewidth}\raggedright
\(R\)
\end{minipage} & \begin{minipage}[b]{\linewidth}\raggedright
\(P\)
\end{minipage} & \begin{minipage}[b]{\linewidth}\raggedright
\(S\)
\end{minipage} & \begin{minipage}[b]{\linewidth}\raggedright
\(\delta^*\)
\end{minipage} \\
\midrule\noalign{}
\endhead
\bottomrule\noalign{}
\endlastfoot
No repair cost & 0 & 0.15 & 1.0 & 85.00 & 48 & 25 & 0 & 0.617 \\
Low repair & 1 & 0.15 & 1.4 & 82.75 & 48 & 15 & −3.75 & 0.513 \\
\textbf{Baseline} & 2 & 0.15 & 1.8 & 80.50 & 48 & 5 & −7.50 &
\textbf{0.430} \\
High repair & 4 & 0.15 & 2.6 & 76.00 & 48 & −15 & −15.00 & 0.308 \\
Very high & 6 & 0.15 & 3.4 & 71.50 & 48 & −35 & −22.50 & 0.250 \\
Cheap exploit & 2 & 0.05 & 1.8 & 93.50 & 48 & 5 & −2.50 & 0.514 \\
Expensive exploit & 2 & 0.30 & 1.8 & 61.00 & 48 & 5 & −15.00 & 0.232 \\
\end{longtable}

\begin{tablenote}

The friction-regime rows above vary the repair multiplier \(\kappa\)
(the physical knob from the Second Law: see Theorem 6 \citep{TC-IV})
holding \(\psi^{\text{off}} = 0.15\), \(\psi^{\text{def}} = 0.25\),
\(\theta = 0.15\), \(c = 2\), \(E_R = 100\). Because \(\Phi\), \(T\),
and \(S\) all vary with \(\kappa\), they are recomputed consistently for
each row. The exploitation-cost rows hold \(\kappa = 2\) and vary
\(\theta\).

\emph{Note on game type.} The PD/Chicken boundary is
\(\Phi_{\text{PD}} = 2 + 4\theta\kappa\psi^{\text{def}}\). The
very-high-repair row (\(\kappa = 6\) with
\(\Phi = 3.4 > \Phi_{\text{PD}} = 2.9\)) is a Chicken game
(\(T > R > S > P\)), for which the grim-trigger threshold is
\(\delta^* = (T - R)/(T - S)\), not \((T - R)/(T - P)\), because the
deviator's best reply to perpetual punishment is \(C\) once \(S > P\)
(Theorem~\ref{thm-cooperation-theorem}). The high-repair row
(\(\kappa = 4\)) sits exactly on the boundary
\(\Phi = \Phi_{\text{PD}} = 2.6\), where \(P = S = -15\) and the two
threshold forms coincide (\(\delta^* = 0.308\)). All remaining rows,
including both exploitation-cost variants, are Prisoner's Dilemmas
(\(\Phi < \Phi_{\text{PD}}\)), computed with
\(\delta^* = (T - R)/(T - P)\). The friction rows fall monotonically in
\(\kappa\) (\(0.513, 0.430, 0.308, 0.250\) for \(\kappa = 1, 2, 4, 6\)),
driven by the punishment payoffs \(P\) and \(S\) worsening as \(\kappa\)
rises. At the calibration baseline \(\kappa = 2\) and above, these lie
strictly below the canonical \(0.5\); at the baseline \(\theta\) and
damage fractions the threshold falls below \(0.5\) for all
\(\kappa \gtrsim 1.2\), so the comparison is knife-edge only near the
admissibility boundary \(\kappa \to 1^+\) (recall that \citep{TC-IV}
derives \(\kappa > 1\) for any physical system, calibrates repair
energetics to roughly \(1.2\)--\(1.7\), and takes \(\kappa = 2\) as a
representative baseline at the high end of that range). The excluded
\(\kappa = 0\) comparison row reaches \(0.617\) and the boundary
\(\kappa = 1\) row \(0.513\), reflecting that with little or no repair
surcharge the punishment for mutual defection is far milder and
cooperation is correspondingly harder. Very cheap exploitation
(\(\theta = 0.05\)) likewise lifts \(\delta^*\) slightly above \(0.5\)
(to \(0.514\)) even at \(\kappa = 2\); every physically admissible
regime tabulated here nonetheless remains below \(0.6\). One further
caveat: in the PD regime the grim punishment path repeats the stage Nash
equilibrium \((D,D)\), so grim trigger is subgame-perfect there; in the
Chicken regime the grim punishment profile \((D,D)\) is not a stage Nash
equilibrium (with \(S > P\), \(D\) is not a best reply to \(D\)), so
grim trigger sustains cooperation as a Nash equilibrium only.

\end{tablenote}

\textbf{Key findings:}

\begin{enumerate}
\def\labelenumi{\arabic{enumi}.}
\tightlist
\item
  \textbf{Higher friction → lower \(\delta^*\) → easier cooperation.}
  When conflict is more destructive, the ``punishment'' for mutual
  defection is harsher, making cooperation sustainable even for
  short-sighted agents.
\item
  \textbf{Higher exploitation cost → lower \(\delta^*\).} When defection
  is expensive (well-defended targets), the temptation premium shrinks
  and cooperation is easier.
\item
  \textbf{Across all physically admissible regimes tabulated above,
  \(\delta^* < 0.6\).} Cooperation is sustainable for any agent that
  values future interactions at more than \textasciitilde60\% of present
  ones, a modest requirement. This bound is a property of the tabulated
  regimes, not a theorem of the model: near the admissibility boundary
  \(\kappa \to 1^+\) combined with very cheap exploitation, \(\delta^*\)
  can slightly exceed \(0.6\) (e.g., \(\kappa = 1.05\),
  \(\theta = 0.01\) at the baseline damage fractions gives
  \(\delta^* = 0.603\)).
\end{enumerate}

\section{Network Effects: The Friction
Amplifier}\label{network-effects-the-friction-amplifier}

\subsection{Incorporating Network Friction into
Payoffs}\label{incorporating-network-friction-into-payoffs}

The single-interaction payoffs above omit a crucial physical cost:
defection injects friction into the network (Definition 9
\citep{TC-IV}), degrading all future cooperative transactions for every
agent in the network. This makes defection costlier than the
single-round analysis suggests, and the additional cost enters the
repeated game through the punishment branch of the equilibrium condition
rather than the single-round payoff matrix.

\subsection{The Full Defection Cost}\label{the-full-defection-cost}

Over the infinite horizon, the total cost of a single defection of
severity \(v\) includes the cascading friction (Theorem 10
\citep{TC-IV}):

\[C_{\text{cascade}} = \frac{\eta \cdot v \cdot M \cdot \bar{E}}{\epsilon}\]

This cost is distributed across all \(M\) transactions and all agents in
the network. The fraction borne by the defecting agent (through degraded
future cooperative payoffs) is:

\[C_{\text{defector}}^{\text{cascade}} = \frac{C_{\text{cascade}}}{N} = \frac{\eta \cdot v \cdot M \cdot \bar{E}}{N \cdot \epsilon}\]

\begin{theorem}[Network-Adjusted Cooperation Condition]
\label{thm-network-adjusted-cooperation}

When network friction effects are included, the effective punishment payoff becomes:

$$\tilde{P} = P - \frac{C_{\text{defector}}^{\text{cascade}}}{(1/\epsilon)} = P - \frac{\eta \cdot v \cdot M \cdot \bar{E}}{N}$$

(amortized over the recovery time horizon). The network-adjusted critical discount factor is:

$$\tilde{\delta}^* = \frac{T - R}{T - \tilde{P}} < \delta^*$$

Network friction makes cooperation strictly easier to sustain.
\end{theorem}

\begin{proof}
Since the cascading friction is a real energy cost borne by the defector (via degraded future cooperative opportunities), and $C_{\text{defector}}^{\text{cascade}} > 0$, we have $\tilde{P} < P$. Therefore $T - \tilde{P} > T - P$, and:

$$\tilde{\delta}^* = \frac{T - R}{T - \tilde{P}} < \frac{T - R}{T - P} = \delta^*$$
\end{proof}

\subsection{Numerical Illustration}\label{numerical-illustration}

For the violation severity we set \(v = E_R\), the per-interaction
resource pool. This represents a full-violation baseline: a defector who
appropriates the counterpart's entire resource share causes harm equal
to \(E_R\). This is the natural unit for energy-denominated violations,
and the bound \(v \leq E_R\) follows from resource conservation (a
single interaction cannot extract more than \(E_R\) from the system).
The numerical results scale linearly with \(v\); smaller violations
reduce \(C_{\text{cascade}}\) proportionally.

Using baseline parameters (\(E_R = 100\), \(M = 1000\),
\(\bar{E} = 10\), \(\eta = 0.01\), \(v = E_R = 100\), \(N = 100\),
\(\epsilon = 0.05\)):

\[C_{\text{defector}}^{\text{cascade}} = \frac{0.01 \times 100 \times 1000 \times 10}{100 \times 0.05} = \frac{10{,}000}{5} = 2{,}000\]

Amortized over the recovery period (\(\sim 1/\epsilon = 20\) periods):

\[\tilde{P}_{\text{per period}} \approx P - 2000/20 = 5 - 100 = -95\]

\[\tilde{\delta}^* = \frac{80.5 - 48}{80.5 - (-95)} = \frac{32.5}{175.5} = 0.185\]

\textbf{With network friction, the cooperation threshold drops to
\(\tilde{\delta}^* = 0.185\).} Even extremely short-sighted agents find
cooperation rational when the network punishes defection through
friction cascades.

\subsection{Deception Penalties}\label{deception-penalties}

Defection via strategic signal corruption rather than overt conflict
incurs the micro-friction costs derived in \citep{TC-II}:

\begin{itemize}
\tightlist
\item
  Decision cost to the victim: \(\Delta U_B(q)\) (Theorem 1
  \citep{TC-II})
\item
  Verification overhead on the network: \(\Delta W(q)\) (Theorem 2
  \citep{TC-II})
\item
  Deception-induced friction:
  \(\Delta\phi_{\text{info}} = \eta_{\text{info}} \cdot \Delta H \cdot W_{\text{bit}}\)
  (Definition 6 \citep{TC-II})
\end{itemize}

These costs further increase the effective punishment and reduce
\(\tilde{\delta}^*\). The deception channel does not create a
``cheaper'' form of defection; it incurs different but equally real
thermodynamic costs.

\section{\texorpdfstring{Extension to \(N\)-Player
Games}{Extension to N-Player Games}}\label{extension-to-n-player-games}

\subsection{\texorpdfstring{The \(N\)-Player Public Goods
Game}{The N-Player Public Goods Game}}\label{the-n-player-public-goods-game}

To extend beyond two players, consider the canonical \(N\)-player
interaction: a \textbf{public goods game} with energy-denominated
payoffs.

\begin{definition}[$N$-Player Energy Public Goods Game]
\label{def-n-player-public-goods}

$N$ agents each choose whether to Cooperate ($C$: contribute energy $c_i > 0$ to a shared pool) or Defect ($D$: contribute nothing). The shared pool generates a return multiplied by factor $\alpha > 1$ (the efficiency gain from cooperation) and is distributed equally. Agent $i$'s payoff:

$$\pi_i = \frac{\alpha \sum_{j=1}^{N} c_j \cdot \mathds{1}[j \text{ cooperates}]}{N} - c_i \cdot \mathds{1}[i \text{ cooperates}] - \phi \cdot \bar{E} \cdot \frac{N - n_C}{N}$$

where $n_C$ is the number of cooperators and the last term captures the friction tax from defectors (proportional to the defection rate).
\end{definition}

In our framework, \(\alpha\) is the \textbf{cooperative
multiplier}\footnote{Within this paper, \(\alpha\) denotes the
  cooperative multiplier exclusively. Other papers in the series reuse
  the symbol for unrelated quantities, such as the marginal energy-yield
  coefficient \(\alpha_{ij}\) in the quadratic utility of \citep{TC-I};
  the two should not be conflated when the papers are read together.},
the ratio of total output under cooperation to total input. That
\(\alpha > 1\) is an assumption of the model, motivated qualitatively by
the welfare-dominance result of \citep{TC-I} (Theorem 9 \citep{TC-I},
Corollary 10 \citep{TC-I}): the variational equilibrium allocates
resources more efficiently than the unconstrained regime, but no
quantitative value of \(\alpha\) is derived there. Its specific value
(here \(\alpha = 2\), cooperative doubling) is an illustrative baseline,
not a derived constant.

More concretely, when all \(N\) agents cooperate, each contributes \(c\)
and receives \(\alpha c\) (net gain \((\alpha - 1)c\)). When \(n_C < N\)
agents cooperate, the cooperators pay the coordination cost while
defectors free-ride, but the system suffers friction loss.

\subsection{\texorpdfstring{The One-Shot \(N\)-Player
Equilibrium}{The One-Shot N-Player Equilibrium}}\label{the-one-shot-n-player-equilibrium}

In the single-round game, an agent's switch from \(D\) to \(C\) gains it
\(\alpha c/N\) from the enlarged pool and, because \(n_C\) rises by one,
reduces its own friction tax by \(\phi_0 \bar{E}/N\), at contribution
cost \(c\). Defection is therefore dominant when:

\[\frac{\alpha c}{N} + \frac{\phi_0 \bar{E}}{N} < c \iff N > \alpha + \frac{\phi_0 \bar{E}}{c}\]

That is, defection dominates when the cooperative return to each
individual, plus the avoided share of the friction tax, is less than the
individual's contribution, the standard free-rider problem sharpened by
friction. At the baseline calibration below (\(\alpha = 2\),
\(\phi_0\bar{E} = 5\), \(c = 10\)) the condition is \(N > 2.5\): for all
group sizes \(N \geq 3\) defection is dominant, while at \(N = 2\)
cooperation is one-shot dominant (consistent with \(\delta_2^* = 0\) in
the table below). For large \(N\), defection dominates unless \(\alpha\)
is very large.

For \(N > \alpha + \phi_0\bar{E}/c\), the one-shot NE is universal
defection, yielding payoff \(\pi_i^{DD} = -\phi_0 \bar{E}\) (only the
friction tax from the non-cooperative environment). The cooperative
payoff is \(\pi_i^{CC} = (\alpha - 1)c\). The efficiency loss is:

\[\Delta\pi = (\alpha - 1)c + \phi_0 \bar{E} > 0\]

\subsection{\texorpdfstring{The Repeated \(N\)-Player
Game}{The Repeated N-Player Game}}\label{the-repeated-n-player-game}

\begin{theorem}[N-Player Cooperation Sustainability]
\label{thm-n-player-cooperation}

In the infinitely repeated $N$-player energy public goods game with discount factor $\delta$ and grim trigger punishment (all agents defect forever upon any defection), cooperation is a NE if:

$$\delta \geq \delta_N^* = \frac{\pi_i^{\text{deviate}} - \pi_i^{CC}}{\pi_i^{\text{deviate}} - \pi_i^{DD}}$$

where $\pi_i^{\text{deviate}}$ is agent $i$'s payoff from unilaterally defecting while all others cooperate:

$$\pi_i^{\text{deviate}} = \frac{\alpha (N-1) c}{N} - 0 = \frac{\alpha(N-1)c}{N}$$

(receives the public good without contributing). The critical discount factor is:

$$\delta_N^* = \frac{\frac{\alpha(N-1)c}{N} - (\alpha - 1)c}{\frac{\alpha(N-1)c}{N} + \phi_0 \bar{E}} = \frac{c(N - \alpha)}{(N-1)\alpha c + N\phi_0\bar{E}}$$
\end{theorem}

\begin{proof}
The deviation payoff is $\pi_i^{\text{deviate}} = \alpha(N-1)c/N$ (receives share of the $(N-1)$ cooperators' contributions without paying). The cooperative payoff is $\pi_i^{CC} = \alpha c - c = (\alpha-1)c$. The punishment payoff is $\pi_i^{DD} = -\phi_0\bar{E}$ (universal defection with friction tax). Applying the grim trigger calculation:

$$\delta_N^* = \frac{\pi_i^{\text{deviate}} - \pi_i^{CC}}{\pi_i^{\text{deviate}} - \pi_i^{DD}} = \frac{\alpha(N-1)c/N - (\alpha-1)c}{\alpha(N-1)c/N + \phi_0\bar{E}}$$

Simplifying the numerator:

$$\frac{\alpha(N-1)c - N(\alpha-1)c}{N} = \frac{c[\alpha N - \alpha - \alpha N + N]}{N} = \frac{c(N - \alpha)}{N}$$

So:

$$\delta_N^* = \frac{c(N - \alpha)/N}{\alpha(N-1)c/N + \phi_0\bar{E}} = \frac{c(N-\alpha)}{\alpha(N-1)c + N\phi_0\bar{E}}$$
\end{proof}

\emph{Remark:} The deviation payoff above omits the friction tax
\(\phi_0\bar{E}/N\) that Definition~\ref{def-n-player-public-goods}
imposes on all agents, including the deviator, in the defection round.
Including this term yields the corrected threshold
\(\delta_N^{*,\text{corr}} = [c(N-\alpha) - \phi_0\bar{E}] / [(N-1)(\alpha c + \phi_0\bar{E})]\),
which is strictly smaller than the stated \(\delta_N^*\). The omission
is therefore conservative: cooperation is achievable for any
\(\delta \geq \delta_N^*\) per the theorem, and in fact for the strictly
smaller \(\delta_N^{*,\text{corr}}\) as well. All qualitative
conclusions (cooperation scales with \(N\), friction lowers the
threshold) are strengthened rather than weakened.

\subsection{\texorpdfstring{Properties of
\(\delta_N^*\)}{Properties of \textbackslash delta\_N\^{}*}}\label{properties-of-delta-n-star}

\begin{corollary}[Cooperation Strengthens with Network Friction]
\label{cor-cooperation-strengthens-friction}

$$\frac{\partial \delta_N^*}{\partial \phi_0} < 0$$

Higher network friction (harsher punishment for the all-defect state) makes cooperation easier to sustain.
\end{corollary}

\begin{proof}
The numerator $c(N - \alpha)$ does not depend on $\phi_0$. The denominator $\alpha(N-1)c + N\phi_0\bar{E}$ is strictly increasing in $\phi_0$. Therefore $\delta_N^*$ is strictly decreasing in $\phi_0$.
\end{proof}

\begin{corollary}[Cooperation Strengthens with Scale Under Friction]
\label{cor-cooperation-strengthens-scale}

When friction is non-negligible ($\phi_0 > 0$), the critical discount factor $\delta_N^*$ remains bounded (does not approach 1) as $N \to \infty$:

$$\lim_{N \to \infty} \delta_N^* = \frac{c}{\alpha c + \phi_0 \bar{E}} < 1$$
\end{corollary}

\begin{proof}
Dividing numerator and denominator by $N$:

$$\delta_N^* = \frac{c(1 - \alpha/N)}{\alpha(1 - 1/N)c + \phi_0\bar{E}} \xrightarrow{N \to \infty} \frac{c}{\alpha c + \phi_0\bar{E}}$$

Since $\phi_0 > 0$: $\alpha c + \phi_0\bar{E} > c$ (for $\alpha > 1$ or $\phi_0 > 0$), so the limit is strictly less than 1.
\end{proof}

\textbf{Physical interpretation:} Without friction (\(\phi_0 = 0\)) the
threshold has the ceiling \(\lim_{N \to \infty} \delta_N^* = 1/\alpha\):
cooperation grows progressively harder as the group expands, and in the
limit of marginally beneficial cooperation (\(\alpha \to 1^+\),
vanishing per-capita surplus) the ceiling approaches \(1\), recovering
the canonical large-group free-rider collapse in which no admissible
discount factor sustains cooperation. Non-zero friction removes this:
the limit becomes \(c/(\alpha c + \phi_0 \bar{E})\), strictly below
\(1\) for every \(\alpha\) and every \(N\), because the friction penalty
on universal defection grows with group size while the per-capita
deviation gain saturates. Cooperation therefore remains sustainable as a
Nash equilibrium in arbitrarily large populations, rather than
collapsing with group size.

\subsection{\texorpdfstring{Numerical Illustration: \(N\)-Player
Cooperation
Thresholds}{Numerical Illustration: N-Player Cooperation Thresholds}}\label{n-player-cooperation-thresholds}

\textbf{Parameters:} \(c = 10\), \(\alpha = 2\) (cooperative doubling),
\(\phi_0 = 0.5\) (moderate background friction; \(\phi_0 > 0\) is an
assumed calibration parameter, not a derived quantity --- it is what
keeps the large-\(N\) limit below 1), \(\bar{E} = 10\).

\begin{longtable}[]{@{}
  >{\raggedright\arraybackslash}p{(\linewidth - 6\tabcolsep) * \real{0.0909}}
  >{\raggedright\arraybackslash}p{(\linewidth - 6\tabcolsep) * \real{0.3636}}
  >{\raggedright\arraybackslash}p{(\linewidth - 6\tabcolsep) * \real{0.3636}}
  >{\raggedright\arraybackslash}p{(\linewidth - 6\tabcolsep) * \real{0.1818}}@{}}
\caption{\(N\)-player cooperation threshold \(\delta_N^*\) versus group
size \(N\).}\tabularnewline
\toprule\noalign{}
\begin{minipage}[b]{\linewidth}\raggedright
\(N\)
\end{minipage} & \begin{minipage}[b]{\linewidth}\raggedright
Deviation gain \(\pi^{\text{dev}} - \pi^{CC}\)
\end{minipage} & \begin{minipage}[b]{\linewidth}\raggedright
Punishment gap \(\pi^{\text{dev}} - \pi^{DD}\)
\end{minipage} & \begin{minipage}[b]{\linewidth}\raggedright
\(\delta_N^*\)
\end{minipage} \\
\midrule\noalign{}
\endfirsthead
\toprule\noalign{}
\begin{minipage}[b]{\linewidth}\raggedright
\(N\)
\end{minipage} & \begin{minipage}[b]{\linewidth}\raggedright
Deviation gain \(\pi^{\text{dev}} - \pi^{CC}\)
\end{minipage} & \begin{minipage}[b]{\linewidth}\raggedright
Punishment gap \(\pi^{\text{dev}} - \pi^{DD}\)
\end{minipage} & \begin{minipage}[b]{\linewidth}\raggedright
\(\delta_N^*\)
\end{minipage} \\
\midrule\noalign{}
\endhead
\bottomrule\noalign{}
\endlastfoot
2 & 0 & 10 + 5 = 15 & 0.000 \\
3 & 3.33 & 13.33 + 5 = 18.33 & 0.182 \\
5 & 6.00 & 16 + 5 = 21 & 0.286 \\
10 & 8.00 & 18 + 5 = 23 & 0.348 \\
50 & 9.60 & 19.6 + 5 = 24.6 & 0.390 \\
100 & 9.80 & 19.8 + 5 = 24.8 & 0.395 \\
\(\infty\) & 10.00 & 20 + 5 = 25 & 0.400 \\
\end{longtable}

\textbf{Key result:} Even as \(N \to \infty\), \(\delta_N^*\) converges
to 0.4, well below 1. Cooperation remains a Nash Equilibrium for any
agent with \(\delta \geq 0.4\). Without friction (\(\phi_0 = 0\)),
\(\delta_N^*\) for \(N = 100\) would be \(9.8/19.8 = 0.495\), and for
\(N \to \infty\) it would approach \(10/20 = 0.5\). Friction provides
the extra margin that makes cooperation robust.

\section{Evolutionary Stability: Defection Cannot
Invade}\label{evolutionary-stability-defection-cannot-invade}

\subsection{The Evolutionary Game-Theoretic
Framework}\label{the-evolutionary-game-theoretic-framework}

Beyond asking whether cooperation is a NE, we ask: \textbf{can defection
invade a cooperative population?} This is the evolutionary stability
question \citep{MaynardSmith1973}.

\begin{definition}[Evolutionary Stable Strategy (ESS)]
\label{def-ess}

A strategy $s^*$ is an ESS if, for any mutant strategy $s' \neq s^*$:

(a) $\pi(s^*, s^*) > \pi(s', s^*)$ (\textbf{strict NE}), or

(b) $\pi(s^*, s^*) = \pi(s', s^*)$ and $\pi(s^*, s') > \pi(s', s')$ (\textbf{stability condition}).
\end{definition}

\subsection{TFT as an Approximate ESS}\label{tft-as-an-approximate-ess}

In the iterated Prisoner's Dilemma with energy payoffs, Tit-for-Tat
(TFT) is a well-known strong performer \citep{Axelrod1984}. While TFT is
not a strict ESS in the technical sense (it does not strictly outperform
other ``nice'' strategies like Always-Cooperate in a homogeneous
population), it satisfies a strong \textbf{invasion barrier} condition:

\begin{theorem}[Defection Invasion Barrier]
\label{thm-defection-invasion-barrier}

In a population of $N$ TFT-playing agents with discount factor $\delta > \delta^*$, a mutant Always-Defect (ALLD) agent earns:

$$V_{\text{ALLD}} = T + \frac{\delta P}{1 - \delta}$$

while a TFT agent in the same population earns:

$$V_{\text{TFT}} = \frac{R}{1 - \delta}$$

The ALLD mutant is strictly dominated ($V_{\text{ALLD}} < V_{\text{TFT}}$) when:

$$\delta > \frac{T - R}{T - P} = \delta^*$$

Under the physical payoff model, this holds for all $\delta > 0.430$ (baseline parameters).
\end{theorem}

\begin{proof}
An ALLD mutant, upon meeting a TFT agent, defects in round 1 (getting $T$) and then faces retaliation: mutual defection from round 2 onward (getting $P$ each round). Its lifetime payoff against any TFT agent is:

$$V_{\text{ALLD}} = T + \delta P + \delta^2 P + \cdots = T + \frac{\delta P}{1 - \delta}$$

A TFT agent meeting other TFT agents cooperates every round:

$$V_{\text{TFT}} = R + \delta R + \delta^2 R + \cdots = \frac{R}{1 - \delta}$$

Then $V_{\text{TFT}} > V_{\text{ALLD}}$:

$$\frac{R}{1 - \delta} > T + \frac{\delta P}{1 - \delta}$$

$$R > (1 - \delta)T + \delta P$$

$$\delta > \frac{T - R}{T - P} = \delta^*$$

which is exactly the cooperation condition from Theorem~\ref{thm-cooperation-theorem}.
\end{proof}

\textbf{Physical interpretation:} A defector entering a cooperative
population gets one ``hit'' (the temptation payoff \(T\)) and then finds
itself in perpetual mutual defection, which under physical payoffs means
a punishment stream far below the cooperative reward \(R\) --- and, in
the destructive regime (\(\kappa\psi > 1\)), net-negative returns
(\(P < 0\)). The defector's strategy is self-defeating: it destroys the
cooperative environment it needs to profit from. In evolutionary terms,
the ``invasion fitness'' of defection is negative in a cooperative
population.

\subsection{The Defector's Dilemma: Extinction Through
Friction}\label{the-defectors-dilemma-extinction-through-friction}

Beyond the direct retaliation mechanism, the defector also faces the
cascading friction cost (§\ref{network-effects-the-friction-amplifier}).
Even if a defector can avoid direct retaliation (e.g., in a large
anonymous population), each defection increases the network friction
\(\phi\), which degrades the returns to \emph{all} agents, including the
defector. In a population where the defector's payoff depends on the
cooperative surplus (because even defectors benefit from a low-friction
network for their non-defection interactions), systematic defection is a
self-undermining strategy.

\begin{corollary}[Defection Self-Extinction in Friction Environments]
\label{cor-defection-self-extinction}

In a network with friction coupling, a defector whose defection rate $\nu_D$ injects enough friction to push $\phi$ above the network collapse threshold (Corollary 11 \citep{TC-IV}) destroys the cooperative surplus from which it profits. In the limit, a population of defectors achieves the minimum network payoff $\pi^{DD} = -\phi_0 \bar{E}$, which is strictly less than zero. Defection is not merely dominated; it is self-destructive.
\end{corollary}

\section{Tit-for-Tat Dynamics in Energy
Terms}\label{tit-for-tat-dynamics-in-energy-terms}

\subsection{Short-Term Defection Profit is
Erasure}\label{short-term-defection-profit-is-erasure}

We now formally show that a defector's short-term energy gain is
\textbf{wiped out} over repeated interactions.

\begin{theorem}[Defection Profit Erasure]
\label{thm-defection-profit-erasure}

An agent that defects once against a TFT population and then returns to cooperation incurs a net lifetime loss whenever:

$$\delta > \frac{T - R}{R - S}$$

(without friction). Under the physical payoff model (baseline parameters), the basic threshold is $\delta = 0.586$. When network friction costs are included ($F = 100$ per period over 20 recovery periods), the threshold drops to $\delta \approx 0.183$.
\end{theorem}

\begin{proof}
Consider an agent that defects in period $t$ (getting $T$), is retaliated against in period $t+1$ (getting $S$), and then both agents return to cooperation from period $t+2$ onward. The deviation costs the agent:

\begin{itemize}
\item Period $t$: gains $T$ instead of $R$ (net gain: $T - R$)
\item Period $t+1$: earns $S$ instead of $R$ (net loss: $R - S$)
\item Periods $t+2, \ldots$: back to $R$ (no change)
\end{itemize}
Present-value net gain from the one-shot deviation:

$$\Delta V = (T - R) - \delta(R - S)$$

The deviation is unprofitable ($\Delta V < 0$) when:

$$\delta > \frac{T - R}{R - S}$$

With baseline parameters: $\delta > (80.5 - 48)/(48 - (-7.5)) = 32.5/55.5 = 0.586$.

However, this analysis ignores the network friction cost. Including the per-period friction penalty amortized over recovery time ($\sim 1/\epsilon = 20$ periods), even a single defection imposes a sustained future cost. From §\ref{numerical-illustration}, the per-defector cascade cost is $C_{\text{defector}}^{\text{cascade}} = 2{,}000$, giving $F = 2{,}000 / 20 = 100$ energy units per period of degraded cooperative returns. When we include this:

$$\Delta V_{\text{adj}} = (T - R) - \delta(R - S) - \delta \sum_{k=1}^{1/\epsilon} \delta^{k-1} F$$

Over the recovery horizon $1/\epsilon = 20$ the discounted friction stream is $\delta\sum_{k=1}^{20}\delta^{k-1}F = F\,\delta(1 - \delta^{20})/(1 - \delta)$, indistinguishable at the relevant small $\delta$ from its infinite-horizon envelope $\delta F/(1 - \delta)$ (since $\delta^{20} \approx 0$). Setting $\Delta V_{\text{adj}} = 0$ with this envelope and clearing the factor $(1 - \delta) > 0$:

$$(T - R)(1 - \delta) - \delta(R - S)(1 - \delta) - \delta F = 0$$

$$\iff (R - S)\,\delta^2 - \bigl[(T - R) + (R - S) + F\bigr]\delta + (T - R) = 0.$$

At the baseline ($T - R = 32.5$, $R - S = 55.5$, $F = 100$) this is $55.5\,\delta^2 - 188\,\delta + 32.5 = 0$, whose admissible root is

$$\delta = \frac{188 - \sqrt{188^2 - 4(55.5)(32.5)}}{2(55.5)} = \frac{188 - \sqrt{28129}}{111} \approx 0.183.$$

Including friction thus lowers the erasure threshold from $0.586$ to $\delta \approx 0.183$: the deviation is unprofitable for any agent with $\delta \gtrsim 0.183$.
\end{proof}

\textbf{Physical interpretation:} The defector's profit
(\(T - R = 32.5\)) is a one-time energy windfall. The retaliation cost
(\(R - S = 55.5\)) comes one period later. The friction cost
(\(F = 100\)/period) compounds over \(\sim 20\) periods. For any agent
that expects to survive more than a few interactions
(\(\delta > 0.183\)), the short-term defection profit is negative on
net. Defection is an energetic ``payday loan'' with punishing interest.

\section{System-Level Optimality: The Welfare
Theorem}\label{system-level-optimality-the-welfare-theorem}

\subsection{The Social Welfare
Function}\label{the-social-welfare-function}

Define the \textbf{system welfare} as the total net energy available to
all agents after interaction costs:

\[SW = \sum_{i=1}^{N} \pi_i - \mathcal{C}_{\text{total}}\]

where \(\mathcal{C}_{\text{total}}\) is the total interaction cost
(Definition 12 \citep{TC-IV} plus Definition 7 \citep{TC-II}).

The unweighted sum embeds an assumption that should be stated
explicitly: adding payoffs across agents presumes a common zero and a
common scale, i.e., full interpersonal comparability. The energy
denomination supplies exactly this --- payoffs are measured in Joules, a
shared physical unit, whereas utils correspond to no physical observable
and are not interpersonally comparable (see the Related Work
discussion). This comparability is a modeling premise imported with the
energy denomination, not a result derived within the series. The welfare
statements that follow, including the uniqueness corollary, are made
relative to it.

\subsection{The Maximum Welfare
Theorem}\label{the-maximum-welfare-theorem}

\begin{theorem}[Cooperation Maximizes System Welfare]
\label{thm-cooperation-maximizes-welfare}

Among all strategy profiles in the repeated $N$-player game, the cooperative profile ($C$ for all agents in all periods) uniquely maximizes the per-period system welfare:

$$SW^{CC} = N(\alpha - 1)c \geq SW^{s} \quad \forall \text{ strategy profiles } s$$

with equality only when $s$ is also universally cooperative.
\end{theorem}

\begin{proof}
Under universal cooperation:
\begin{itemize}
\item Total resource input: $Nc$ (all agents contribute)
\item Total output: $\alpha N c$ (cooperative multiplier)
\item Net welfare: $N(\alpha - 1)c$
\item Interaction cost: $\mathcal{C}_{\text{total}} = 0$ (no conflict, no deception, $\phi = 0$)
\end{itemize}
Under any profile with $n_D > 0$ defectors:
\begin{itemize}
\item Resource input: $(N - n_D)c < Nc$ (defectors do not contribute)
\item Output: $\alpha(N - n_D)c < \alpha Nc$ (reduced by lost contributions)
\item Network friction tax: per Definition~\ref{def-n-player-public-goods}, each of the $N$ agents pays $\phi(s) \bar{E} \cdot n_D / N$, for a system total of $n_D \, \phi(s) \bar{E} > 0$
\item Net welfare: strictly less than $N(\alpha - 1)c$
\end{itemize}
Formally, summing the payoffs of Definition~\ref{def-n-player-public-goods}:

$$SW^{CC} - SW^{s} = \underbrace{n_D (\alpha - 1)c}_{\text{lost contributions}} + \underbrace{n_D \, \phi(s) \bar{E}}_{\text{friction tax}} > 0$$

for any $s$ with $n_D > 0$. In addition, defection raises the ambient per-period transaction tax $\phi(s) \cdot M \cdot \bar{E}$ of Definition 12 \citep{TC-IV} ($M$ transactions of average value $\bar{E}$), which enters $SW$ through $\mathcal{C}_{\text{total}}$ and vanishes only under universal cooperation ($\phi = 0$); this contribution only widens the gap above. Uniqueness follows from the strict inequality above: for any strategy profile with $n_D \geq 1$ defectors, the lost-contributions term $n_D(\alpha - 1)c > 0$ ensures that no defecting profile achieves $SW^{CC}$. (*Note: The Tullock contest dissipation term $D(s)$ applies to arms-race conflict settings (see Proposition 1 \citep{TC-IV}) and does not apply in this public goods free-rider context; uniqueness here rests solely on the lost-contributions and friction-tax terms.*)
\end{proof}

\subsection{The System-Level Nash
Equilibrium}\label{the-system-level-nash-equilibrium}

Combining Theorem~\ref{thm-cooperation-theorem} and
Theorem~\ref{thm-cooperation-maximizes-welfare}:

\begin{corollary}[Cooperation as the Unique Efficient Nash Equilibrium]
\label{cor-cooperation-unique-efficient-ne}

In the infinitely repeated energy-based game with $\delta > \delta^*$:

(a) Universal cooperation is a Nash Equilibrium (Theorem~\ref{thm-cooperation-theorem}).

(b) Universal cooperation uniquely maximizes system welfare (Theorem~\ref{thm-cooperation-maximizes-welfare}).

(c) Universal cooperation minimizes total interaction cost ($\mathcal{C}_{\text{total}} = 0$).

Therefore, the cooperative strategy profile ("cooperation") is the unique Nash Equilibrium that is also Pareto-optimal and system-welfare-maximizing.
\end{corollary}

\textbf{Physical interpretation:} Among all Nash Equilibria of the
repeated energy-denominated game, universal cooperation is the unique
profile that maximizes aggregate welfare and minimizes total
dissipation. This is an equilibrium-selection result, not equilibrium
uniqueness: the Folk Theorem continuum persists --- in the Prisoner's
Dilemma regime, where mutual defection is a stage-game equilibrium
(\(P > S\)), universal defection remains a Nash equilibrium at every
discount factor --- and what the framework supplies is a principled
welfare-dominance criterion that singles out cooperation among these
equilibria.

\section{Discussion}\label{discussion}

\subsection{Physical Interpretation}\label{physical-interpretation}

The payoff matrix derived in §\ref{the-energy-denominated-payoff-model}
is determined entirely by the resource value \(E_R\), the exploitation
efficiency \(\theta\), the friction multiplier \(\Phi\), and the
coordination cost \(c\). The Prisoner's Dilemma ordering
\(T > R > P > S\) (Proposition~\ref{prop-physical-prisoners-dilemma})
holds in the parameter region where exploitation is cheap and friction
is moderate; at higher friction (\(\Phi > \Phi_{\text{PD}}\)) the
ordering becomes that of Hawk-Dove, with \(S > P\), where the
grim-trigger threshold takes the form \((T-R)/(T-S)\) rather than
\((T-R)/(T-P)\). Both orderings emerge from the same physical model
under different parameter values, and the location of the transition
(\(\Phi_{\text{PD}} = 2 + 4\theta\kappa\psi^{\text{def}}\)) is derived
within the model --- from the friction quantities of \citep{TC-IV}
together with the exploitation efficiency \(\theta\) introduced here ---
rather than chosen.

The negativity of \(P\) under \(\Phi > 2\) reflects the Second Law as it
acts on contested boundaries, combined with a magnitude condition on
damage. Boundary structures are organized configurations that required
directed work to assemble; an attack disrupts that organization in a
single irreversible step, while restoration requires sustained work
against the entropic gradient. The condition \(\Phi > 2\) states that
the work cost of repair, amplified by the irreversibility factor
\(\kappa > 1\), exceeds the work value extracted from the resource; it
is equivalent to combined boundary damage exceeding \(1/\kappa\) (one
half at \(\kappa = 2\)). The Second Law forces \(\kappa > 1\) but not
the damage magnitudes, so \(\Phi > 2\) marks the destructive regime
rather than a thermodynamic necessity; at the mild baseline
(\(\kappa\psi < 1\)) mutual defection is costly but net-positive. In the
destructive regime mutual defection is net-dissipative: more energy is
converted to waste heat than is captured as usable resource.

The cooperation threshold \(\delta^*\) admits a corresponding physical
reading. The discount factor \(\delta = \sigma e^{-r}\) encodes
per-period survival probability \(\sigma\) and time-preference rate
\(r\). An agent satisfying \(\delta \geq \delta^*\) is one whose
expected continuation in the system is sufficient that future energy
access dominates the immediate energy windfall of seizure. The threshold
\(\delta^* \approx 0.430\) is the level of expected persistence at which
an agent's own thermodynamic future outweighs the one-shot temptation;
the cooperation theorem (Theorem~\ref{thm-cooperation-theorem}) is the
statement that an agent intending to persist will not defect on a
cooperator.

The friction state variable \(\phi_t\) tracks accumulated entropy in
shared infrastructure. Each defection injects \(\Delta\phi = \eta v\)
into the network, and recovery proceeds at rate \(\epsilon\) per period
\citep{TC-IV}. The repeated-game payoff coupling through \(\phi_t\)
(Theorem~\ref{thm-network-adjusted-cooperation}) is the second-law cost
of restoration distributed across all participants: a defector exports
entropy onto the network, and subsequent cooperative transactions pay
the work cost of importing free energy to drive \(\phi_t\) back toward
baseline. The network-adjusted threshold
\(\tilde{\delta}^* \approx 0.185\) reflects this additional cost;
defection becomes more difficult to sustain because each instance
damages the substrate of all future interactions.

Welfare uniqueness (Theorem~\ref{thm-cooperation-maximizes-welfare},
Corollary~\ref{cor-cooperation-unique-efficient-ne}) follows from the
same dissipation accounting. Universal cooperation is the only profile
in which contest investment \(e^*\), boundary damage, and friction
injection are all zero. Every other profile expends energy on contest,
repair, or network restoration. The welfare-maximizing equilibrium is
therefore the profile that minimizes irreversible work per period, and
the result is not a coincidence of utility specification but a property
of the underlying energy accounting.

\subsection{Testable Predictions}\label{testable-predictions}

\begin{longtable}[]{@{}
  >{\raggedright\arraybackslash}p{(\linewidth - 4\tabcolsep) * \real{0.3333}}
  >{\raggedright\arraybackslash}p{(\linewidth - 4\tabcolsep) * \real{0.2222}}
  >{\raggedright\arraybackslash}p{(\linewidth - 4\tabcolsep) * \real{0.4444}}@{}}
\caption{Testable predictions and their falsification
criteria.}\tabularnewline
\toprule\noalign{}
\begin{minipage}[b]{\linewidth}\raggedright
Prediction
\end{minipage} & \begin{minipage}[b]{\linewidth}\raggedright
Formal source
\end{minipage} & \begin{minipage}[b]{\linewidth}\raggedright
Falsification criterion
\end{minipage} \\
\midrule\noalign{}
\endfirsthead
\toprule\noalign{}
\begin{minipage}[b]{\linewidth}\raggedright
Prediction
\end{minipage} & \begin{minipage}[b]{\linewidth}\raggedright
Formal source
\end{minipage} & \begin{minipage}[b]{\linewidth}\raggedright
Falsification criterion
\end{minipage} \\
\midrule\noalign{}
\endhead
\bottomrule\noalign{}
\endlastfoot
Cooperation threshold \(\delta^* \approx 0.43\) under baseline friction
& Theorem~\ref{thm-cooperation-theorem} & Cooperation thresholds in
energy-denominated experimental games consistently exceed \(0.5\) \\
Single-deviation profit is net-negative for \(\delta > 0.183\) &
Theorem~\ref{thm-defection-profit-erasure} & Single defectors in
iterated games with endogenous reputation costs show net-positive
lifetime payoffs \\
\(N\)-player cooperation threshold converges to \(\sim 0.4\) as
\(N \to \infty\) & Theorem~\ref{thm-n-player-cooperation} & Large-group
experiments with friction show \(\delta_N^* \to 1\) as in the standard
public goods game \\
Network collapse at critical violation frequency &
Corollary~\ref{cor-defection-self-extinction}, Corollary 11
\citep{TC-IV} & Institutional-trust data shows no threshold effect
relating violation frequency to cooperative breakdown \\
Defection invasion fitness is negative in cooperative populations &
Theorem~\ref{thm-defection-invasion-barrier} & Defector mutants
consistently invade and persist in cooperative populations under
physical payoff conditions \\
\end{longtable}

\subsection{Limitations}\label{limitations}

The game-theoretic core inherits standard simplifications from the
repeated-games literature:

\begin{enumerate}
\def\labelenumi{\arabic{enumi}.}
\item
  \textbf{Binary strategy space.} Agents choose \(C\) or \(D\). Real
  interactions admit continuous gradations of compliance. A continuous
  extension, where agents choose a degree of compliance
  \(\lambda \in [0,1]\) (distinct from the cooperative multiplier
  \(\alpha\) used in the \(N\)-player public goods game), would refine
  the payoff landscape but would not eliminate the friction asymmetry
  driving the result. We conjecture that the cooperative equilibrium
  persists and that Corollary~\ref{cor-cooperation-unique-efficient-ne}
  generalizes to a cooperation manifold.
\item
  \textbf{Indefinite horizon.} The cooperation results assume no
  anticipated final period. For agents with long planning horizons, this
  is well-matched; applications to agents with known finite lifespans
  require the finite-horizon extension with reputation effects.
\item
  \textbf{Perfect monitoring.} The grim trigger strategy assumes
  deviations are observed immediately. Under imperfect monitoring, the
  standard result is that cooperation remains a NE but requires a higher
  \(\delta > \delta^*_{\text{noisy}} > \delta^*\)
  \citep{GreenPorter1984, Fudenberg1991}. The model's baseline
  \(\delta^* = 0.430\) provides substantial headroom before cooperation
  ceases to be sustainable.
\end{enumerate}

\textbf{Scope.} The payoff model assumes symmetric agents contesting a
single resource per interaction. Asymmetric agents (different \(E_R\),
\(\theta\), or \(\kappa\)) would produce asymmetric payoff matrices
requiring bargaining-theoretic extensions. The qualitative result
(cooperation as the unique welfare-maximizing equilibrium) is expected
to hold under asymmetry, but the quantitative thresholds would depend on
the specific asymmetry structure.

\subsection{Applications}\label{applications}

The energy-denominated game theory framework applies to domains where
interactions involve measurable resource costs and repeated engagement:

\begin{enumerate}
\def\labelenumi{\arabic{enumi}.}
\item
  \textbf{Interstate conflict and arms races.} The friction multiplier
  \(\Phi\) provides a formal counterpart to the economic cost of
  military conflict between nations, and the cooperation threshold gives
  a quantitative criterion for when the threat of mutual defection
  (deterrence) suffices to sustain a peaceful equilibrium.
\item
  \textbf{Market competition and oligopoly.} Firms competing for market
  share face coordination costs (trade agreements, regulatory
  compliance) and conflict costs (price wars, patent litigation). The
  model predicts that industries with high friction costs
  (capital-intensive sectors, regulated markets) should exhibit more
  stable cooperative behavior (cartels, industry standards) than
  low-friction industries.
\item
  \textbf{Biological cooperation and altruism.} The cooperative
  equilibrium results extend naturally to weighted payoff functions that
  model inclusive fitness, coupled-agent systems, and belief-extended
  discount factors. These applications, which unify Hamilton's rule,
  reciprocal altruism, and group selection under a single
  payoff-expansion framework, are developed in \citep{Ethics-I}.
\end{enumerate}

\subsection{Future Work}\label{future-work}

\textbf{Continuous strategy spaces.} Extending from binary \(\{C, D\}\)
to continuous compliance \(\lambda \in [0,1]\) with energy-denominated
payoffs, to characterize the full cooperation manifold and verify that
partial defection is dominated.

\textbf{Imperfect monitoring.} Incorporating noisy observation channels,
where defection is detected with probability \(p < 1\), to bound
\(\delta^*_{\text{noisy}}\) as a function of detection probability and
determine when the baseline headroom (\(\delta^* = 0.430\) versus the
canonical \(0.5\)) is sufficient.

\textbf{Heterogeneous agents.} Agents with asymmetric endowments, time
preferences, or friction parameters, to determine whether the welfare
theorem generalizes to bargaining solutions and whether the equilibrium
selection result survives asymmetry.

\textbf{Coalition formation.} Modeling intermediate structures between
pairwise interaction and universal cooperation, to characterize when and
how stable coalitions form and how they resist internal defection.

\textbf{Embeddedness-dependent cooperation threshold.} The cooperation
threshold \(\delta^* \approx 0.430\) is derived for agents in a
shared-pool environment. Embedding agents in a production network, where
each agent's output feeds other agents' inputs, would lower this
threshold for agents with high in-degree and out-degree, because
defection would additionally sever supply lines that downstream agents
must rebuild at thermodynamic cost. Such an extension would provide a
first-principles derivation of why specialized, division-of-labor
economies exhibit more stable cooperation than autarkic ones.

\textbf{Endogenizing the cooperative multiplier.} The multiplier
\(\alpha\) is treated here as a parameter, grounded qualitatively in the
welfare-dominance of the variational equilibrium \citep{TC-I} but not
derived to a specific value. Deriving \(\alpha\) from the underlying
production and allocation structure, for example the input-output
couplings of a production network, would replace the illustrative
baseline with a first-principles value and sharpen the scale results,
whose ceiling \(1/\alpha\) depends directly on it.

\section{Conclusion}\label{conclusion}

We have proved that in the repeated multi-agent game with
energy-denominated payoffs, universal cooperation is the unique Nash
Equilibrium that simultaneously satisfies Pareto efficiency and welfare
maximization, for agents with \(\delta \geq \delta^*\)
(Theorem~\ref{thm-cooperation-theorem},
Theorem~\ref{thm-cooperation-maximizes-welfare},
Corollary~\ref{cor-cooperation-unique-efficient-ne}). This is an
equilibrium-selection result: the Folk Theorem's continuum of equilibria
persists, and what the framework supplies is a principled
welfare-dominance criterion that singles out cooperation among them. The
result rests on the structural properties of the payoff matrix, derived
from thermodynamic parameters rather than stipulated, which make mutual
defection net-negative in the destructive damage regime
(\(\kappa\psi > 1\)) and cooperation self-reinforcing through endogenous
friction dynamics. Crucially, the load-bearing \(N\)-player
welfare-uniqueness result rests on \(\Phi\)-independent machinery ---
the threshold \(\delta_N^*\) carries no friction multiplier --- so it is
robust to the calibration of the damage parameters; the two-player
regime map supplies the mechanism, not the conclusion.

In the two-player game the cooperation threshold is below the canonical
\(0.5\) of the standard Prisoner's Dilemma, and drops further to
\(\tilde{\delta}^* \approx 0.185\) when network friction is included
(Theorem~\ref{thm-network-adjusted-cooperation}). Defection is
evolutionarily unstable (Theorem~\ref{thm-defection-invasion-barrier}),
self-extinguishing (Corollary~\ref{cor-defection-self-extinction}), and
unprofitable even from a single deviation
(Theorem~\ref{thm-defection-profit-erasure}). The result extends to
arbitrarily large populations: without friction the threshold rises
toward the ceiling \(1/\alpha\), approaching \(1\) as cooperation
becomes marginally beneficial and recovering the canonical large-group
free-rider collapse, whereas non-zero background friction holds it
strictly below \(1\) for every group size
(Corollary~\ref{cor-cooperation-strengthens-scale}).

\emph{Necessary Constraints} \citep{TC-I}, \emph{Strategic Entropy
Injection} \citep{TC-II}, and \emph{Thermodynamic Friction}
\citep{TC-IV} provide the physical infrastructure from which the payoff
structure derives, and \emph{Value Dynamics} \citep{TC-VI} analyzes the
dynamical stability of the resulting cooperative configurations.

\appendix

\section{\texorpdfstring{Appendix: Recalled Results from the
\emph{Thermodynamics of Cooperation}
Series}{Appendix: Recalled Results from the Thermodynamics of Cooperation Series}}\label{appendix-recalled-results-from-the-thermodynamics-of-cooperation-series}

This appendix collects, for the reader's convenience, the results that
\emph{Cooperative Equilibrium} imports from earlier papers in the
\emph{Thermodynamics of Cooperation} series: \emph{Necessary
Constraints} \citep{TC-I}, \emph{Strategic Entropy Injection}
\citep{TC-II}, and \emph{Thermodynamic Friction} \citep{TC-IV}. Each
statement is recalled without proof; the proof appears in the cited
source. Nothing here is original to the present paper, and nothing in
the main text depends on this appendix beyond the results stated below.

\subsection{Setup and the cooperative allocation (from Necessary
Constraints)}\label{setup-and-the-cooperative-allocation-from-necessary-constraints}

We consider \(N \geq 2\) self-maintaining agents that persist by
importing free energy and exporting entropy while sharing a
finite-resource environment. All utilities, costs, and payoffs are
denominated in energy units (Joules), following the ecological-economics
convention \citep{Lotka1922, GeorgescuRoegen1971, GeorgescuRoegen1979}.
Each agent \(i\) has a utility \(U_i\) satisfying the regularity
assumption of \citep{TC-I} (Assumption 1 \citep{TC-I}:
twice-differentiable, strictly concave, monotone-increasing at the
origin, coercive) and a boundary of integrity \(B_i > 0\) in energy
units.

\textbf{Variational equilibrium and the cooperative split} (Theorem 7
\citep{TC-I}). When all agents optimize simultaneously under shared
resource constraints, the solution concept is the variational
(Normalized Nash) equilibrium, at which every agent faces the same
per-unit shadow price on each scarce resource. For symmetric agents
contesting a resource of energy value \(E_R > 0\), the equilibrium
allocates the resource equally. Net of a coordination cost
\(c \ll E_R\), this equal split is the cooperative payoff used in the
main text, \(R = E_R/2 - c\).

\textbf{Welfare dominance} (Theorem 9 \citep{TC-I}, Corollary 10
\citep{TC-I}). The variational equilibrium uniquely maximizes total
welfare over the feasible set, strictly dominating greedy, equal-split,
and every other feasible allocation. This is the basis for the
cooperative multiplier exceeding unity in the \(N\)-player public goods
game.

\subsection{Conflict costs and the friction multiplier (from
Thermodynamic
Friction)}\label{conflict-costs-and-the-friction-multiplier-from-thermodynamic-friction}

When an agent violates a boundary constraint, the resulting contest and
boundary damage are governed by the following quantities recalled from
\citep{TC-IV}: the offensive and defensive damage fractions
\(\psi^{\text{off}}, \psi^{\text{def}} \in [0,1]\), the boundary damage
inflicted on attacker and victim respectively (with combined boundary
damage \(\psi = \psi^{\text{off}} + \psi^{\text{def}}\)); the repair
multiplier \(\kappa \geq 1\), with \(\kappa > 1\) whenever the Second
Law applies, because reconstructing a damaged boundary requires both
entropy clearing and negentropy import where maintenance requires only
the latter; and the friction multiplier \(\Phi = 1 + \kappa\psi\). The
main text additionally uses the exploitation efficiency
\(\theta \in (0,1)\), the fraction of a resource's value an attacker
must spend to seize it from an unprepared cooperator; this quantity is
defined in the present paper, not recalled from the series.

\textbf{Symmetric contest} (Proposition 1 \citep{TC-IV}). In the
symmetric two-agent Tullock contest over \(E_R\), each agent invests
\(e^* = E_R/4\) at the Nash equilibrium, and total contest dissipation
is \(E_R/2\).

\textbf{Net-negative conflict} (Theorem 6 \citep{TC-IV}). Mutual
defection dissipates more energy than the contested resource is worth
whenever \(\Phi\) exceeds the population-dependent threshold
\(N/(N-1)\); for two agents this is \(\Phi > 2\). The condition
\(\Phi > 2\) is equivalent to \(\kappa\psi > 1\), i.e.~combined boundary
damage exceeding \(1/\kappa\) (one half at \(\kappa = 2\));
\(\kappa > 1\) and strictly positive damage alone do not force it. For
three or more contestants the threshold falls to \(\Phi > N/(N-1)\),
equivalently \(\kappa\psi > 1/(N-1)\) --- a threshold that shrinks as
contestants are added and vanishes only in the large-\(N\) limit, where
any positive damage suffices. This result is what drives the
mutual-defection payoff net-negative in the destructive regime of the
main text.

\subsection{Friction dynamics (from Thermodynamic
Friction)}\label{friction-dynamics-from-thermodynamic-friction}

Boundary violations degrade not only the immediate interaction but all
subsequent cooperative transactions, through a network friction
coefficient \(\phi \geq 0\) that taxes per-transaction overhead
(monitoring, verification, precautionary defense). A violation of
severity \(v\) injects \(\Delta\phi = \eta v\) into \(\phi\), where
\(\eta > 0\) is the friction sensitivity; absent new violations,
\(\phi\) decays geometrically at rate \(\epsilon \in (0,1)\) per period;
and the combined law of motion is
\(\phi_{t+1} = (1-\epsilon)\phi_t + \eta v_t \cdot \mathds{1}_{\{\text{violation at } t\}}\).
The injection is additive and instantaneous while recovery is
multiplicative and gradual, the thermodynamic asymmetry between
increasing and decreasing entropy.

\textbf{Cascading friction} (Theorem 10 \citep{TC-IV}). A single
violation of severity \(v\) produces total cascading cost
\(\eta v M \bar{E} / \epsilon\) over the full recovery horizon, where
\(M\) is the number of cooperative transactions per period and
\(\bar{E}\) their average value; the amplification ratio
\(\eta M \bar{E}/\epsilon\) spans \(10^3\) to \(10^4\) for modest
network parameters. In the main text this cost enters the repeated game
through the punishment branch of the equilibrium condition.

\subsection{The deception channel (from Strategic Entropy
Injection)}\label{the-deception-channel-from-strategic-entropy-injection}

Defection may also take the form of strategic entropy injection into a
communication channel \citep{TC-II} rather than overt contest. This
carries its own strictly positive energy costs: a decision cost on the
deceived receiver (Theorem 1 \citep{TC-II}), a verification-energy
burden on a defender (Theorem 2 \citep{TC-II}), and a mapped
contribution to the friction coefficient \(\phi\) (Definition 6
\citep{TC-II}). These costs enter the punishment payoff analogously to
physical friction and are recalled in the main text only at the point of
use.

\addcontentsline{toc}{section}{References}
\begin{thebibliography}{28}
\providecommand{\natexlab}[1]{#1}
\providecommand{\url}[1]{\texttt{#1}}
\expandafter\ifx\csname urlstyle\endcsname\relax
  \providecommand{\doi}[1]{doi: #1}\else
  \providecommand{\doi}{doi: \begingroup \urlstyle{rm}\Url}\fi

\bibitem[Fudenberg and Maskin(1986)]{Fudenberg1986}
Drew Fudenberg and Eric Maskin.
\newblock The folk theorem in repeated games with discounting or with
  incomplete information.
\newblock \emph{Econometrica}, 54\penalty0 (3):\penalty0 533--554, 1986.
\newblock \doi{10.2307/1911307}.

\bibitem[Fudenberg and Tirole(1991)]{Fudenberg1991}
D~Fudenberg and J~Tirole.
\newblock Game theory, 616 pages, 1991.

\bibitem[Lostracco(2026{\natexlab{a}})]{TC-I}
Keith Lostracco.
\newblock Thermodynamics of cooperation: Necessary constraints,
  2026{\natexlab{a}}.
\newblock URL \url{https://doi.org/10.5281/zenodo.19635522}.
\newblock \doi{10.5281/zenodo.19635522}.

\bibitem[Lostracco(2026{\natexlab{b}})]{TC-IV}
Keith Lostracco.
\newblock Thermodynamics of cooperation: Thermodynamic friction,
  2026{\natexlab{b}}.
\newblock URL \url{https://doi.org/10.5281/zenodo.19635849}.
\newblock \doi{10.5281/zenodo.19635849}.

\bibitem[Lostracco(2026{\natexlab{c}})]{TC-II}
Keith Lostracco.
\newblock Thermodynamics of cooperation: Strategic entropy injection,
  2026{\natexlab{c}}.
\newblock URL \url{https://doi.org/10.5281/zenodo.19635813}.
\newblock \doi{10.5281/zenodo.19635813}.

\bibitem[Von~Neumann and Morgenstern(1944)]{VonNeumann1944}
John Von~Neumann and Oskar Morgenstern.
\newblock \emph{Theory of Games and Economic Behavior}.
\newblock Princeton University Press, 1944.

\bibitem[Nash(1950)]{Nash1950}
John~F. Nash.
\newblock Equilibrium points in {N}-person games.
\newblock \emph{Proceedings of the National Academy of Sciences}, 36\penalty0
  (1):\penalty0 48--49, 1950.
\newblock \doi{10.1073/pnas.36.1.48}.

\bibitem[Nash(1951)]{Nash1951}
John~F. Nash.
\newblock Non-cooperative games.
\newblock \emph{Annals of Mathematics}, 54\penalty0 (2):\penalty0 286--295,
  1951.
\newblock \doi{10.2307/1969529}.

\bibitem[Maynard~Smith and Price(1973)]{MaynardSmith1973}
John Maynard~Smith and George~R. Price.
\newblock The logic of animal conflict.
\newblock \emph{Nature}, 246\penalty0 (5427):\penalty0 15--18, 1973.
\newblock \doi{10.1038/246015a0}.

\bibitem[Hamilton(1964)]{Hamilton1964a}
William~D. Hamilton.
\newblock The genetical evolution of social behaviour. {I}.
\newblock \emph{Journal of Theoretical Biology}, 7\penalty0 (1):\penalty0
  1--16, 1964.
\newblock \doi{10.1016/0022-5193(64)90038-4}.

\bibitem[Axelrod(1984)]{Axelrod1984}
Robert Axelrod.
\newblock \emph{The Evolution of Cooperation}.
\newblock Basic Books, 1984.

\bibitem[Taylor and Jonker(1978)]{Taylor1978}
Peter~D. Taylor and Leo~B. Jonker.
\newblock Evolutionarily stable strategies and game dynamics.
\newblock \emph{Mathematical Biosciences}, 40\penalty0 (1--2):\penalty0
  145--156, 1978.
\newblock \doi{10.1016/0025-5564(78)90077-9}.

\bibitem[Weibull(1995)]{Weibull1995}
J{\"o}rgen~W. Weibull.
\newblock \emph{Evolutionary Game Theory}.
\newblock MIT Press, 1995.

\bibitem[Osborne and Rubinstein(1994)]{OsborneRubinstein1994}
Martin~J. Osborne and Ariel Rubinstein.
\newblock \emph{A Course in Game Theory}.
\newblock MIT Press, 1994.

\bibitem[Szab{\'o} and F{\'a}th(2007)]{SzaboFath2007}
Gy{\"o}rgy Szab{\'o} and G{\'a}bor F{\'a}th.
\newblock Evolutionary games on graphs.
\newblock \emph{Physics Reports}, 446\penalty0 (4--6):\penalty0 97--216, 2007.
\newblock \doi{10.1016/j.physrep.2007.04.004}.

\bibitem[Perc et~al.(2013)Perc, G{\'o}mez-Garde{\~n}es, Szolnoki, Flor{\'i}a,
  and Moreno]{PercEtAl2013}
Matja{\v z} Perc, Jes{\'u}s G{\'o}mez-Garde{\~n}es, Attila Szolnoki, Luis~M.
  Flor{\'i}a, and Yamir Moreno.
\newblock Evolutionary dynamics of group interactions on structured
  populations: a review.
\newblock \emph{Journal of the Royal Society Interface}, 10\penalty0
  (80):\penalty0 20120997, 2013.
\newblock \doi{10.1098/rsif.2012.0997}.

\bibitem[Perc et~al.(2017)Perc, Jordan, Rand, Wang, Boccaletti, and
  Szolnoki]{PercEtAl2017}
Matja{\v z} Perc, Jillian~J. Jordan, David~G. Rand, Zhen Wang, Stefano
  Boccaletti, and Attila Szolnoki.
\newblock Statistical physics of human cooperation.
\newblock \emph{Physics Reports}, 687:\penalty0 1--51, 2017.
\newblock \doi{10.1016/j.physrep.2017.05.004}.

\bibitem[Santos et~al.(2008)Santos, Santos, and
  Pacheco]{SantosSantosPacheco2008}
Francisco~C. Santos, Marta~D. Santos, and Jorge~M. Pacheco.
\newblock Social diversity promotes the emergence of cooperation in public
  goods games.
\newblock \emph{Nature}, 454\penalty0 (7201):\penalty0 213--216, 2008.
\newblock \doi{10.1038/nature06940}.

\bibitem[Jusup et~al.(2022)Jusup, Holme, Kanazawa, Takayasu, Romi{\'c}, Wang,
  Ge{\v c}ek, Lipi{\'c}, Podobnik, Wang, Luo, Klanj{\v s}{\v c}ek, Fan,
  Boccaletti, and Perc]{JusupEtAl2022}
Marko Jusup, Petter Holme, Kiyoshi Kanazawa, Misako Takayasu, Ivan Romi{\'c},
  Zhen Wang, Sun{\v c}ana Ge{\v c}ek, Tomislav Lipi{\'c}, Boris Podobnik, Lin
  Wang, Wei Luo, Tin Klanj{\v s}{\v c}ek, Jingfang Fan, Stefano Boccaletti, and
  Matja{\v z} Perc.
\newblock Social physics.
\newblock \emph{Physics Reports}, 948:\penalty0 1--148, 2022.
\newblock \doi{10.1016/j.physrep.2021.10.005}.

\bibitem[Anttila and Annila(2011)]{AnttilaAnnila2011}
Jani Anttila and Arto Annila.
\newblock Natural games.
\newblock \emph{Physics Letters A}, 375\penalty0 (43):\penalty0 3755--3761,
  2011.
\newblock \doi{10.1016/j.physleta.2011.08.056}.

\bibitem[Adami and Hintze(2018)]{AdamiHintze2018}
Christoph Adami and Arend Hintze.
\newblock Thermodynamics of evolutionary games.
\newblock \emph{Physical Review E}, 97\penalty0 (6):\penalty0 062136, 2018.
\newblock \doi{10.1103/PhysRevE.97.062136}.

\bibitem[Lotka(1922)]{Lotka1922}
Alfred~J. Lotka.
\newblock Contribution to the energetics of evolution.
\newblock \emph{Proceedings of the National Academy of Sciences}, 8\penalty0
  (6):\penalty0 147--151, 1922.
\newblock \doi{10.1073/pnas.8.6.147}.

\bibitem[Georgescu-Roegen(1971)]{GeorgescuRoegen1971}
Nicholas Georgescu-Roegen.
\newblock \emph{The Entropy Law and the Economic Process}.
\newblock Harvard University Press, 1971.

\bibitem[Hirshleifer(1991)]{Hirshleifer1991}
Jack Hirshleifer.
\newblock The technology of conflict as an economic activity.
\newblock \emph{American Economic Review}, 81\penalty0 (2):\penalty0 130--134,
  1991.

\bibitem[Green and Porter(1984)]{GreenPorter1984}
Edward~J. Green and Robert~H. Porter.
\newblock Noncooperative collusion under imperfect price information.
\newblock \emph{Econometrica}, 52\penalty0 (1):\penalty0 87--100, 1984.
\newblock ISSN 00129682, 14680262.
\newblock URL \url{http://www.jstor.org/stable/1911462}.

\bibitem[Lostracco(2026{\natexlab{d}})]{Ethics-I}
Keith Lostracco.
\newblock A theory of ethics from first principles: Thermodynamic norms,
  2026{\natexlab{d}}.
\newblock URL \url{https://doi.org/10.5281/zenodo.19635922}.
\newblock \doi{10.5281/zenodo.19635922}.

\bibitem[Lostracco(2026{\natexlab{e}})]{TC-VI}
Keith Lostracco.
\newblock Thermodynamics of cooperation: Value dynamics, 2026{\natexlab{e}}.
\newblock URL \url{https://doi.org/10.5281/zenodo.19635864}.
\newblock \doi{10.5281/zenodo.19635864}.

\bibitem[Georgescu-Roegen(1979)]{GeorgescuRoegen1979}
Nicholas Georgescu-Roegen.
\newblock Comments on the papers by {Daly} and {Stiglitz}.
\newblock In V.~Kerry Smith, editor, \emph{Scarcity and Growth Reconsidered},
  pages 95--105. Johns Hopkins University Press, Baltimore, 1979.

\end{thebibliography}
\end{document}
